%%%%%%%%%%%%%%%%%%%%%%%%%%%%%%%%%%%%%%%%%%%%%%%%%%%%%%
%% "Magic comments" for the LaTeX editor -------------
% !TEX root = ../Diss.tex
%%%%%%%%%%%%%%%%%%%%%%%%%%%%%%%%%%%%%%%%%%%%%%%%%%%%%%
%
\begin{showExamples}%
%
%%%%%%%%%%%%%%%%%%%%%%%%%%%%%%%%%%%%%%%%%%%%%%%%%%%%%%%%%%%%
\chapter{Anleitung zur Nutzung \mbox{dieser Vorlage}}%
\label{chap:Examples}
%%%%%%%%%%%%%%%%%%%%%%%%%%%%%%%%%%%%%%%%%%%%%%%%%%%%%%%%%%%%
%
%
Dieses Kapitel zeigt den grundlegenden Aufbau eines einfachen Kapitels.
Die einzelnen Abschnitte beschreiben die Struktur der Vorlage und geben wichtige allgemeine Tipps.
Die Verwendung einzelner Features dieser Vorlage wird in \cref{chap:Examples} detailliert beschrieben und demonstriert.

Es wird empfohlen, die einzelnen Unterkapitel jedes Kapitels 
als eigene Dateien anzulegen und sie mit dem \verb+\input{}+-Befehl einzubinden
(s. Quelltext der Datei \printfilepathquoted{content/Inhalt-Examples.tex}).
Dies erlaubt die einzelnen Unterkapitel bei Bedarf leichter zu verschieben 
oder mit einem einzigen \%-Zeichen temporär auszukommentieren 
und erleichtert so die Fehlersuche.

Teile des Manuskriptes lassen sich ebenfalls ausblenden, ohne sie auskommentieren zu müssen,
indem man den auszublendenden Text zwischen
\lstinline|\begin{ShowInhaltXXX}| und \lstinline|\end{ShowInhaltXXX}| einschließt.
Hierzu ist in der Datei \printfilepathquoted{preambel/AlleSchalter.tex}
eine entsprechende \LaTeX{}"=Umgebung mittels
\lstinline|\setUserDefinedBoolean{showInhaltXXX}{Bool-Wert}| zu definieren und
auf \printkeyword{true} (Inhalt einblenden) 
oder \printkeyword{false} (Inhalt ausblenden) zu setzen.%
%
\footnote{Die Befehle
\lstinline|\setUserDefinedBoolean{IrgendeinName}{Bool-Wert}| 
sind spezielle Makros, die jeweils eine Boolean-Variable mit dem entsprechenden Wert definieren,
sowie eine \LaTeX"=Umgebung \printkeyword{IrgendeinName},
Alles, was im weiteren \LaTeX-Quellcode zwischen
\lstinline|\begin{IrgendeinName}| und \lstinline|\end{IrgendeinName}| steht,
wird je nach dem spezifizierten Wert entweder angezeigt (\printkeyword{true}) 
oder ausgeblendet (\printkeyword{false}).
Der Vorteil dieser Vorgehensweise gegenüber einem einfachen Auskommentieren 
liegt \ua darin, dass die zwischen
\lstinline|\begin{IrgendeinName}| und \lstinline|\end{IrgendeinName}|
eingebundenen Dateien weiterhin im Verzeichnisbaum vom TeXnicCenter \etc
sichtbar bleiben und bei Bedarf schnell als Referenz aufgerufen werden können.}
%
%
%% Unterkapiteln mit speziellen Hinweisen und Beispielen:
%
%%%%%%%%%%%%%%%%%%%%%%%%%%%%%%%%%%%%%%%%%%%%%%%%%%%%%%%%%%%%
\section{Kompilierung der Vorlage}%
\label{sec:Kompilierung}
%%%%%%%%%%%%%%%%%%%%%%%%%%%%%%%%%%%%%%%%%%%%%%%%%%%%%%%%%%%%
%
Um diese Vorlage nutzen zu können, benötigt man eine \LaTeX-Distribution (\zB \printswname{MiKTeX}\footnote{\url{http://www.miktex.org/}} oder \printswname{TeXLive}\footnote{\url{https://tug.org/texlive/acquire.html}}).
Sofern man nicht die riesengroße Komplettinstallation wählt, wird beim ersten Kompilieren eine Internet"=Verbindung benötigt, um Zusatz"=\glspl{gls:pkg} dynamisch nachladen zu können.
Zur Erstellung des Glossars und des Abkürzungsverzeichnisses wird zusätzlich Perl benötigt.
Unter Windows müsste hierzu zusätzlich beispielsweise \printswname{ActivePerl}\footnote{\url{https://www.activestate.com/products/activeperl/}}
oder \printswname{StrawberryPerl}\footnote{\url{http://strawberryperl.com/}} installiert werden.
Bei den Linux-Distributionen ist Perl automatisch mit dabei.

Zur bequemen Bearbeitung der \LaTeX"=Quellcode"=Dateien empfiehlt sich die Verwendung einer guten \LaTeX"=Entwicklungsumgebung in Kombination mit einem geeigneten, sprich SyncTeX"=fähigen, PDF"=Betrachter.
Manche Entwicklungsumgebungen bieten eine integrierte PDF"=Vorschau, welche Vor- und Nachteile haben kann.
Ein wichtiges Kriterium bei der Wahl der Entwicklungsumgebung ist die Möglichkeit, zwischen den einzelnen Stellen im Quellcode und im PDF hin und her springen zu können.
Unter Windows war lange Zeit \textbf{TeXnicCenter}\footnote{\url{http://www.texniccenter.org/}} in Kombination mit \printswname{SumatraPDF}\footnote{\url{https://www.sumatrapdfreader.org}} eine gute Wahl gewesen.
Mittlerweile tendieren die meisten dazu, \printswname{TeXstudio}\footnote{\url{https://www.texstudio.org/}} zu verwenden.
Dieses bietet eine eingebaute PDF"=Vorschau und viele nützliche Features und ist sowohl unter Windows als auch unter Linux verfügbar.

%%%%%%%%%%%%%%%%%%%%%%%%%%%%%%%%%%%%%%%%%%%%%%%%%%%%%%%%%%%%
\subsection{MiKTeX-Einstellungen}
\label{sec:MiKTeX}
%%%%%%%%%%%%%%%%%%%%%%%%%%%%%%%%%%%%%%%%%%%%%%%%%%%%%%%%%%%%
Sofern MiKTeX als \LaTeX"=Distribution verwendet wird, sollte man darauf achten, dass bei Bedarf Zusatzpakete vom Internet dynamisch nachgeladen werden können.
Sofern diese Option nicht bei der Installation gesetzt worden ist, kann sie nachträglich in der MiKTeX-Console aktiviert werden.
Zu finden ist diese im Startmenü unter \enquote{MiKTeX}, \enquote{MiKTeX Console (Admin)}.
Unter \enquote{Settings} findet sich ein Reiter \enquote{General}, wo im Bereich \enquote{Package installation} entweder die Option \enquote{Always install missing packages on the fly} oder \enquote{Ask me} ausgewählt werden soll.
Sofern sich der Rechner im Netzwerk des Fraunhofer IOSB befindet, muss zusätzlich noch die Proxy-Option korrekt gesetzt werden.
Hierfür muss man im Bereich \enquote{Package installation} der MiKTeX Console auf \enquote{Change} gehen und bei \enquote{Connection Settings} die Verwendung des Proxy-Servers \printkeyword{proxy-ka.iosb.fraunhofer.de} mit Port \printkeyword{80} aktivieren (s. \cref{fig:MiKTeX-Proxy}).
Wird dies nicht gemacht, können benötigte Pakete nicht nachgeladen werden.

\begin{figure}[H]%
\centering%
\includegraphics[width=\linewidth]{./images/examples/MiKTeX-Proxy.png}%
\caption{MiKTeX-Einstellungen zum Nachladen der Zusatzpakete}%
\label{fig:MiKTeX-Proxy}%
\end{figure}

Nach der MiKTeX-Installation sollte man im Startmenü gleich die MiKTeX-Console aufrufen, den Proxy eintragen und das Update durchführen, um die aktuellste Version der vorinstallierten Pakete zu erhalten.
So vermeidet man beim automatischen Nachladen weiterer \glspl{gls:pkg} aus dem \acrshort{ac:CTAN}-Repository etwaige Inkompatibilitäten aufgrund veralteter Stammpakete.


%%%%%%%%%%%%%%%%%%%%%%%%%%%%%%%%%%%%%%%%%%%%%%%%%%%%%%%%%%%%
\subsection{TeXLive-Einstellungen}
\label{sec:TeXLive}
%%%%%%%%%%%%%%%%%%%%%%%%%%%%%%%%%%%%%%%%%%%%%%%%%%%%%%%%%%%%
Bei Verwendung von TeXLive unter Linux muss darauf geachtet werden, dass alle notwendigen Linux-\glspl{gls:pkg} installiert sind.
Unter Ubuntu 17.04 sollte es funktionieren, wenn folgende Linux-\glspl{gls:pkg} installiert wurden:
\begin{itemize*}
	\item \printkeyword{texlive}
	\item \printkeyword{texlive-lang-german}
	\item \printkeyword{texlive-fonts-extra}
	\item \printkeyword{texlive-bibtex-extra}
	\item \printkeyword{fonts-linuxlibertine}
	\item \printkeyword{biber}
	\item \printkeyword{xindy}
	\item \printkeyword{texmaker}
\end{itemize*}
Texmaker ist eine IDE für \LaTeX, die aber vermutlich über Dependencies schon einige Pakete mitbringt.


%%%%%%%%%%%%%%%%%%%%%%%%%%%%%%%%%%%%%%%%%%%%%%%%%%%%%%%%%%%%
\subsection{Kompilieraufrufe}
\label{sec:Aufruf}
%%%%%%%%%%%%%%%%%%%%%%%%%%%%%%%%%%%%%%%%%%%%%%%%%%%%%%%%%%%%
Zur erfolgreichen Kompilierung des Dokumentes müssen mehrere
Kommandozeilenprogramme aufgerufen werden.
Bei Verwendung einer integrierten Entwicklungsumgebung (IDE),
wird diese so konfiguriert, dass die Aufrufe aus der IDE heraus erfolgen und
die etwaigen Warnungen, Erfolgs- und Fehlermeldungen in der IDE sichtbar werden.
Die entsprechenden Einstellungen für TeXnicCenter und TeXstudio finden sich
in den nachfolgenden Kapiteln.

Die einzelnen Aufrufe sind:
\begin{itemize*}
\item \index{xelatex}\printkeyword{xelatex} zur eigentlichen Kompilierung von \LaTeX-Quelltext,
\item \index{biber}\printkeyword{biber} zur Kompilierung von Bibliografien,
\item \index{makeglossaries}\printkeyword{makeglossaries}, welches intern \index{xindy}\printkeyword{xindy} aufruft,
zur Erstellung einer Zwischenausgabe für das Abkürzungsverzeichnis und das Glossar
\index{makexindex}\printkeyword{makeindex}, welches durch die Verwendung des Pakets \pkg{imakeidx} implizit aufgerufen wird, zur Erstellung des Stichwortverzeichnisses
\end{itemize*}
Bei einem Durchlauf von \printkeyword{xelatex}, \printkeyword{biber}, \printkeyword{makeglossaries} und \printkeyword{makeindex}
werden die einzelnen Inhalte sowie die entsprechenden Querverweise
zuerst jeweils in eine oder mehrere Zwischendateien hinausgeschrieben,
die sodann wieder eingelesen und verarbeitet werden müssen.
Manche Inhalte werden daher erst jeweils beim zweiten Aufruf von \printkeyword{xelatex} generiert.
Für die korrekte Generierung eines Dokumentes mit allen Verzeichnissen
(Inhalts-, Abbildungs-, Tabellen-, Literatur, Abkürzungs-, Begriffs- und Stichwortverzeichnis),
PDF"=Lesezeichen und korrekt gesetzten Querverweisen,
muss \printkeyword{xelatex} daher mehrmals aufgerufen werden.

Sollte man ausnahmsweise das Dokument doch noch aus der Kommandozeile oder von einem Script heraus aufrufen wollen,
so ist die Reihenfolge der Aufrufe wie folgt:
\begin{verbatim}
#> xelatex -synctex=1 -interaction=nonstopmode Diss.tex
#> biber Diss
#> makeglossaries Diss
#> xelatex -synctex=1 -interaction=nonstopmode Diss.tex
#> xelatex -synctex=1 -interaction=nonstopmode Diss.tex
#> xelatex -synctex=1 -interaction=nonstopmode Diss.tex
\end{verbatim}

Wenn kein Kompilierfehler aufgetreten ist,
sollte nach dem vierten Durchlauf von \printkeyword{xelatex} die
\index{Warnung!Please re-run latex}Warnung \enquote{Please re-run latex}
verschwunden sein.
%%%%%%%%%%%%%%%%%%%%%%%%%%%%%%%%%%%%%%%%%%%%%%%%%%%%%%%%%%%%
\section{Speichererweiterung}%
\index{Speichererweiterung}%
\label{sec:Speichererweiterung}
%%%%%%%%%%%%%%%%%%%%%%%%%%%%%%%%%%%%%%%%%%%%%%%%%%%%%%%%%%%%
%
Sollte bei der Kompilierung die \index{Kompilierfehler!TeX capacity exceeded}\index{Fehler!TeX capacity exceeded}Fehler-Meldung \enquote{\texttt{TeX capacity exceeded}} auftreten,
bedeutet dies, dass die von der \LaTeX-Distribution vorgesehene Arbeitsspeicherkapazität
nicht ausreicht, um die Datenmenge zu verarbeiten.
Wenn im Code kein Fehler vorliegt (\zB eine vergessene Klammer,
die ebenfalls für eine solche Fehlermeldung sorgen kann),
kann dies auch daran liegen, dass die Verarbeitung umfangreicher TiKZ-Bilder oder Bibliographien 
Zusatzspeicher benötigen.
Um den Speicher zu erweitern, gibt es zwei Möglichkeiten.
Zum einen kann man bei jedem Aufruf von \texttt{xelatex} die Optionen zur Speichererweiterung,
\zB \printkeyword{-main-memory=500000000}, \printkeyword{-extra-mem-top=500000000}, \printkeyword{-extra-mem-bot=500000000} und ggf. weitere setzen.
Zum anderen kann man diese Einstellungen ein für alle Male direkt bei den LaTeX-Distributionen setzen.

Unter TeXLive kann die Einstellung in der Datei
\printfilepathquoted{/Pfad-zur-TeXLive/texmf.cnf} vorgenommen werden.
Der Aufruf erfolgt am besten, indem man im Terminal den Befehl \printkeyword{kpsewhich -a texmf.cnf} eintippt.
In der Datei kann dann beispielsweise folgendes gesetzt werden:
\begin{latex}[caption={[Einstellungen zur erweiterten Speichernutzung]Einstellungen zur erweiterten Speichernutzung in der Datei \printfilepathquoted{texmf.cnf} bei TeXLive bzw. \printfilepathquoted{xelatex.ini} bei MiKTeX},label={lst:Speichererweiterung}]
main_memory=5000000
extra_mem_top=5000000
extra_mem_bot=5000000
pool_size=5000000
buf_size=5000000
save_size=79999
\end{latex}
Der maximal mögliche Wert für \printkeyword{main\_memory} \etc ist \printkeyword{79999999}.

Unter MiKTeX unterscheidet sich die Vorgehensweise je nach dem,
ob man die Einstellung nur für den aktuellen Benutzer oder
global für alle Benutzer des Computers machen möchte.
Im zweiten Fall braucht man Admin-Rechte.

Um die Einstellungen nur für den aktuellen Benutzer zu setzen,
öffnet man die Windows-Eingabeaufforderung
(Auf den \printkeyword{Start}-Knopf gehen und \printkeyword{cmd} eintippen)
und tippt dort \printkeyword{initexmf --edit-config-file=xelatex} ein.
Um die Einstellungen für alle Benutzer zu setzen,
öffnet man die Windows"=Eingabeaufforderung im Administrator"=Modus
(dazu auf den Menüeintrag \printkeyword{Eingabeaufforderung} mit der rechten Maustaste anklicken und \printkeyword{Als Administrator ausführen} wählen)
und tippt dort \printkeyword{initexmf \textbf{--admin} --edit-config-file=xelatex} ein.

Dabei öffnet sich die Datei \printfilepathquoted{xelatex.ini}, in welcher die \og Einstellungen (s. \cref{lst:Speichererweiterung}) gesetzt werden sollten.%
\footnote{Die Datei \printfilepathquoted{xelatex.ini} mit benutzerübergreifend gültigen Einstellungen befindet sich unter
\printfilepathquoted{C:\bs Program Files\bs MiKTeX\bs miktex\bs config}.
Die dort gesetzten Einstellungen gelten, sofern der einzelne Benutzer keine eigenen Einstellungen definiert hat.
Die benutzerspezifische Datei, deren Einstellungen die globalen Einstellungen überrufen können, befindet sich unter
\printfilepathquoted{C:\bs Users\bs <Benutzername>\bs AppData\bs Roaming\bs MiKTeX\bs 2.9\bs miktex\bs config}.}

Nach Speicherung der Datei ist bei der benutzerspezifischen Anpassung in der Eingabeaufforderung der Befehl
\printkeyword{initexmf --dump=xelatex} aufzurufen.
Im Falle einer benutzerübergreifenden Anpassung ist der Befehl
\printkeyword{initexmf \textbf{--admin} --dump=xelatex}
aufzurufen.
%%%%%%%%%%%%%%%%%%%%%%%%%%%%%%%%%%%%%%%%%%%%%%%%%%%%%%%%%%%%
\section{Aufbau der Vorlage}%
\label{sec:AufbauDerVorlage}
%%%%%%%%%%%%%%%%%%%%%%%%%%%%%%%%%%%%%%%%%%%%%%%%%%%%%%%%%%%%
%
Die Vorlage besteht aus mehreren Dateien und Verzeichnissen.
Ihre Bedeutung ist in \cref{tab:StrukturDerVorlage} zusammengefasst.

%%% Übersicht über die Dateien und Verzechnisse der Vorlage %%%
%%
%% Durch Definition einer Gruppe mit \begingroup ... \endgroup 
%% sicherstellen, dass Anpassungen nur lokal wirken
%% und sich nicht auf die nachfolgenden Tabellen erstrecken
\begingroup%
%% kleinere Schrift
\small%
% Abstand zwischen den Zeilen auf das 2fache setzen
\renewcommand{\arraystretch}{2.0}%
% Seitliche Abstände links und rechts eliminieren
%\setlength\LTleft{0pt}%
%\setlength\LTright{0pt}%
\setlength{\tabcolsep}{0pt}%
\setlength{\topskip}{0pt}%
%\setlength{\LTpre}{0pt}%
% Tabelle aus einer Extra-Datei einlesen und auf die Seitenbreite strecken
\LTXtable{\columnwidth}{tables/DateienUndVerzeichnisse.tex}%
\endgroup%


%%%%%%%%%%%%%%%%%%%%%%%%%%%%%%%%%%%%%%%%%%%%%%%%%%%%%%%%%%%%
\subsection[Projektdatei \enquote{Diss.tcp}]{Projektdatei \enquote{Diss.tcp}}%
\index{Projektdatei}%
\label{sec:Projektdatei}
%%%%%%%%%%%%%%%%%%%%%%%%%%%%%%%%%%%%%%%%%%%%%%%%%%%%%%%%%%%%
%
Bei der Datei \printfilepathquoted{Diss.tcp} handelt es sich um die Projektdatei für den \LaTeX-Editor TeXnicCenter.
In ihr werden die projektbezogenen Einstellungen des TeXnicCenter festgehalten.
Das sind \ua Angaben zur Hauptdatei des Projektes und zur Projektsprache.
Die korrekte Angabe der Projektsprache ist insofern wichtig, als dass diese in TeXnicCenter ab Version 2.0 Beta 1 zur Bestimmung der Sprache für die Rechtschreibprüfung verwendet wird.
Die entsprechenden Einstellungen können im TeXnicCenter über den Menüeintrag \printkeyword{Projekt} $\rightarrow$ \printkeyword{Eigenschaften} vorgenommen werden.


%%%%%%%%%%%%%%%%%%%%%%%%%%%%%%%%%%%%%%%%%%%%%%%%%%%%%%%%%%%%
\subsection[Hauptdatei \enquote{Diss.tex}]{Hauptdatei \enquote{Diss.tex}}%
\index{Hauptdatei}%
\label{sec:Hauptdatei}
%%%%%%%%%%%%%%%%%%%%%%%%%%%%%%%%%%%%%%%%%%%%%%%%%%%%%%%%%%%%
%
Die Hauptdatei ist die Datei \printfilepathquoted{Diss.tex}.
Sie ist verhältnismäßig kurz, da die Hauptinhalte in andere Dateien ausgelagert sind, welche mit Hilfe des \lc{input\{\}} \bzw des \lc{include\{\}}-Befehls eingebunden werden.
Die Hauptdatei besteht im Wesentlichen aus drei Abschnitten.
Im ersten stehen die sogenannten \index{Magic Comments}\enquote{Magic comments}, mit deren Hilfe manche \LaTeX-IDEs (wie \zB TeXStudio) sich selbst vorkonfigurieren können.
Sie fangen mit \enquote{\texttt{\%~!TeX}} an und geben an, welche Kodierung für die Dateien verwendet wird und welche Programme für die Kompilierung des Quelltextes und der Bibliografie verwendet werden sollen.
Außerdem kann damit angegeben werden, welche Sprache für die Rechtschreibprüfung innerhalb der IDE verwendet werden soll.
Da es aktuell keine dokumentierte Möglichkeit gibt, Aufruf für \printkeyword{makeglossaries} miteinzubinden,
führt eine Übernahme der durch die \enquote{Magic Comments} vordefinierten Reihenfolge dazu, dass keine Glossare generiert werden.
Daher sind die \enquote{Magic Comments} aktuell durch ein zusätzliches Prozentzeichen deaktiviert.

Im zweiten Abschnitt wird die Header-Datei eingebunden.
In dieser wird die verwendete Dokumentklasse (inklusive Papierformat und Schriftgröße) definiert sowie weitere Dateien eingebunden.
In diesen werden die zu landenden Pakete, Layout"=Parameter sowie alle weiteren Einstellungen und Makros definiert und konfiguriert.

Im dritten Abschnitt werden die einzelnen Inhalte der Arbeit eingebunden.

Die einzelnen \index{Kapitel}Kapitel der Arbeit werden im Verzeichnis \printfilepathquoted{./content/} als separate Dateien gespeichert.
Es empfiehlt sich als Dateiname das Schema \printfilepathquoted{nn-name.tex} zu verwenden, wobei \printkeyword{nn} die Nummer des Kapitels ist,
sodass die Dateien in der semantisch richtigen Reihenfolge sortiert angezeigt werden.
Die einzelnen Dateien werden per \verb+include{}+-Direktive in der Datei \printfilepathquoted{Diss.tex} eingebunden.
Theoretisch wäre es an dieser Stelle auch möglich mit \verb+\input{}+ zu arbeiten, was jedoch seine Nachteile hätte.
Der Unterschied zwischen den beiden Befehlen wird \href{https://texwelt.de/wissen/fragen/32/was-ist-der-unterschied-zwischen-include-and-input}{auf texwelt.de} erklärt:

\begin{quote}
{\small
\verb+\input{file}+ lädt die Datei an Ort und Stelle in die Ziel-Datei und ist äquivalent
als ob man den Text in \printkeyword{file} direkt in die Ziel-Datei geschrieben hätte.
\verb+\input+ kann letztlich überall für jede Art Datei verwendet werden und kann auch verschachtelt angewendet werden,
d.h. eine eingebundene Datei kann ihrerseits Dateien mit \verb+\input+ einbinden.

\verb+\include{file}+ hingegen führt zunächst einmal ein \verb+\clearpage+ aus bevor es \verb+\input{file}+ ausführt.
Im Gegensatz zu \verb+\input+ kann eine Datei, die mit \verb+\include+ eingebunden wird,
kein weiteres \verb+\include+ enthalten, es ist also keine verschachtelte Anwendung möglich.
Eine mit \verb+\include+ eingebundene Datei kann aber natürlich \verb+\input+ enthalten.
\verb+\include+ erzeugt eine neue \printfilepathquoted{aux}-Datei für die eingebundene Datei.
Das erlaubt es beispielsweise, ein Dokument in mehrere logische Einheiten zu zerlegen (etwa einzelne Kapitel),
die jede einer Datei entsprechen, die mit \verb+\include+ in die Hauptdatei eingebunden wird.
\verb+\includeonly{file1,file3}+ würde dann erlauben, nur gerade bearbeitete Dateien für die Kompilation einzubinden
und durch die separaten \printfilepathquoted{aux}-Dateien dennoch korrekte Seitenzahlen und Querverweise zu erhalten.
Es gibt auch das \pkg{excludeonly} Paket, dessen Befehl \verb+\excludeonly+ das gegensätzliche Verhalten bietet.%
}%
\footnote{\url{https://texwelt.de/wissen/fragen/32/was-ist-der-unterschied-zwischen-include-and-input}}
\end{quote}

Es empfiehlt sich, die bestehende Struktur (zumindest jedoch die beiden Dateien
\printfilepathquoted{Inhalt-FrontMatter.tex} und \printfilepathquoted{Inhalt-Backmatter.tex})
als Vorlage zu übernehmen.

Das entstehende PDF heißt genauso wie die Hauptdatei.

Zur besseren Übersicht und zur Vereinfachung der Fehlersuche wird empfohlen,
die einzelnen Unterkapitel ebenfalls als separate Dateien in Unterverzeichnissen von \printfilepathquoted{./content/} anzulegen
und sie mit den \verb+\input{}+-Direktiven in die jeweiligen Kapitel-Dateien einzubinden.



%%%%%%%%%%%%%%%%%%%%%%%%%%%%%%%%%%%%%%%%%%%%%%%%%%%%%%%%%%%%
\subsection[Bibliografie-Dateien]{Bibliografie-Dateien}%
\index{Bibliografie}%
\label{sec:Bibliografie}
%%%%%%%%%%%%%%%%%%%%%%%%%%%%%%%%%%%%%%%%%%%%%%%%%%%%%%%%%%%%
%
Es wird davon ausgegangen, dass sich sämtliche Bibliografie-Angaben in der Datei
\printfilepathquoted{./bib/Diss.bib} befinden.
Sollten mehrere Bibliografie-Dateien verwendet werden, können diese in der Datei
\printfilepathquoted{./preamble/AllePfade.tex} gesetzt werden.
Details dazu gibt es im \cref{sec:Literaturverzeichnisse}.

%%%%%%%%%%%%%%%%%%%%%%%%%%%%%%%%%%%%%%%%%%%%%%%%%%%%%%%%%%%%
\subsection[Bilder und Zeichnungen]{Bilder und Zeichnungen}%
\index{Bild!Binär-}%
\index{Bild!Vektor-}%
\label{sec:BilderUndZeichnungen}
%%%%%%%%%%%%%%%%%%%%%%%%%%%%%%%%%%%%%%%%%%%%%%%%%%%%%%%%%%%%
%
Bilder \bzw \index{Zeichnung|see{Bild}}Zeichnungen werden auf zwei Arten eingebunden.
Bilder im \index{Bild!Binär-}Binärformat (PNG, JPEG, TIFF, PDF, \etc)
werden mit \lc{includegraphics}-Befehl eingebunden. 
Bei den \index{Bild!TikZ}\gls{gls:tikz}-Zeichnungen handelt es sich um reguläre TeX-Quellcode-Dateien,
die mit dem \verb+\input+-Befehl eingebunden werden.
Für eine einfache Verwaltung wird empfohlen, Binärbilder im Verzeichnis \printfilepathquoted{./images/} abzulegen.
Zusätzliche Pfade können in der Datei \printfilepathquoted{./preambel/AllePfade.tex} definiert werden.
Die \gls{gls:tikz}-Quellcode-Dateien sollten im Verzeichnis \printfilepathquoted{./figures-src/} abgelegt werden.
Während des Kompilierens werden für jede \gls{gls:tikz}-Zeichnung im Verzeichnis \printfilepathquoted{./figures-compiled/} mehrere Dateien erzeugt.
Der Inhalt dieses Verzeichnisses kann gefahrlos gelöscht werden.
Weitere Hinweise und Beispiele zur Einbindung von Grafiken finden sich in \cref{sec:Bilder}.

%%%%%%%%%%%%%%%%%%%%%%%%%%%%%%%%%%%%%%%%%%%%%%%%%%%%%%%%%%%%
\subsection[Datei \enquote{AlleAngaben.tex}: Angaben zur Arbeit]{Datei \enquote{AlleAngaben.tex}: Angaben zur Arbeit}%
\index{Autor}%
\index{Titel}%
\label{sec:Angaben}
%%%%%%%%%%%%%%%%%%%%%%%%%%%%%%%%%%%%%%%%%%%%%%%%%%%%%%%%%%%%
%
Die wichtigsten Einstellungen, die auf jeden Fall geändert werden müssen,
finden sich in der Datei \printfilepathquoted{./preambel/AlleAngaben.tex}.
Hier werden \ua Angaben zum Verfasser, Art und Titel der Arbeit sowie zu den Gutachtern gemacht.

%%%%%%%%%%%%%%%%%%%%%%%%%%%%%%%%%%%%%%%%%%%%%%%%%%%%%%%%%%%%
\subsection[Datei \enquote{AlleSchalter.tex}: Anzeigeeinstellungen]{Datei \enquote{AlleSchalter.tex}: Anzeigeeinstellungen}%
\index{Schalter}%
\label{sec:Schalter}
%%%%%%%%%%%%%%%%%%%%%%%%%%%%%%%%%%%%%%%%%%%%%%%%%%%%%%%%%%%%
%
Eine weitere wichtige Datei ist \printfilepathquoted{./preambel/AlleSchalter.tex}.
Darin wird \ua die \index{Sprache!Hauptsprache}Hauptsprache des Manuskriptes gesetzt, was sich an mehreren Stellen auswirkt.
So wird beispielsweise bei \index{Sprache!globale Umstellung}Umstellung auf Englisch als Hauptsprache
\enquote{Danksagung} durch \enquote{Acknowledgments},
\enquote{Inhaltsverzeichnis} durch \enquote{Contents}
\usw ersetzt.
Auch die Regeln der \index{Silbentrennung!allgemeine Einstellung}Silbentrennung werden entsprechend angepasst.

Durch das Umsetzen der Schalter 
\lstinline|\setUserDefinedBoolean{SchalterName}{Bool-Wert}| 
können hier Teile der Arbeit aus- und wieder eingeblendet werden, ohne dass sie auskommentiert werden müssen.
Neben diesen Einstellungen befinden sich in der Datei \printfilepathquoted{./preambel/AlleSchalter.tex} auch Einstellungen,
die für die Vorbereitung des Manuskriptes zum Druck über den \gls{ac:KSP} wichtig sind.
Insbesondere können hier

\begin{itemize} % Aufzählungen mit mehr als drei Zeilen sollten im Blocksatz stehen, evtl. lässt sich das mit einer Enumerate-Umgebung realisieren

\item \justifying mit der Einstellung
\lstinline|\setUserDefinedBoolean{showFrame}{true}|
die Satzspiegel"=Ränder angezeigt werden und so kontrolliert werden,
ob nichts hinausragt und ob Tabellen und Bilder die komplett verfügbare Breite ausfüllen,

\item \justifying mit den Einstellungen
\lstinline|\setUserDefinedBoolean{coloredlistings}{false}| und 
\lstinline|\setUserDefinedBoolean{coloredlinks}{false}|
dafür gesorgt werden, dass alle Quellcode"=Listings sowie Querverweise und URL"=Adressen,
die normalerweise farbig sind, für die Druckversion nicht farbig gesetzt werden
und so die Anzahl der farbig zu druckenden Seiten reduziert wird,

\item \justifying mit der Einstellung
\lstinline|\setUserDefinedBoolean{useCMYKcolors}{true}|
eine Farbkonvertierung aller Farben in den CMYK-Farbraum für den Offset-Druck vorgenommen werden.
Diese Einstellung sollte nur für die Druckversion vorgenommen werden.%
\footnote{Da manche RGB-Farben bei der Konvertierung in den CMYK-Farbraum blass aussehen,
ist es sinnvoll, eine Alternativversion der betroffenen Farben im CYMK-Farbraum zu definieren,
die gut aussieht.
Beispiele dafür gibt es in der Datei \printfilepathquoted{./preambel/ColorSettings.tex}.
Vorzugsweise sollten aber die KIT-Corporate-Identity-Farben verwendet werden,
welche in der Datei \printfilepathquoted{KAcolors.sty} definiert sind.
Für diese Farben wurde sowohl eine RGB- als auch eine CMYK-Definition erstellt.}

\end{itemize}

%%%%%%%%%%%%%%%%%%%%%%%%%%%%%%%%%%%%%%%%%%%%%%%%%%%%%%%%%%%%
\subsection[Datei \enquote{Hyphenation.tex}: Silbentrennung von Spezialwörtern]{Datei \enquote{Hyphenation.tex}: \mbox{Silbentrennung von Spezialwörtern}}%
\index{Silbentrennung}%
\label{sec:Hyphenation}
%%%%%%%%%%%%%%%%%%%%%%%%%%%%%%%%%%%%%%%%%%%%%%%%%%%%%%%%%%%%
%
Regeln zur Silbentrennung unbekannter Wörter (\zB Fachbegriffe)
können in der Datei \printfilepathquoted{./preambel/Hyphenation.tex} festgelegt werden.
Zu beachten ist, dass zusammengesetzte Wörter mit einem Bindestrich
ausschließlich am Bindestrich getrennt werden,
wogegen auch ein Eintrag in die Datei \printfilepathquoted{Hyphenation.tex} nicht hilft.
Um Silbentrennung an anderen Stellen eines zusammengesetzten Wortes zu erlauben,
muss man bei Verwendung zusammengesetzter Wörter mit Bindestrich diesen durch „\verb+"=+“ ersetzen.
Alternative Trennstellen in bestimmten Wörtern kann man mit Hilfe der Hinweise auf 
\href{https://de.wikibooks.org/wiki/LaTeX-W%C3%B6rterbuch:_Silbentrennung}{wikibooks.org}%
\footnote{\url{https://de.wikibooks.org/wiki/LaTeX-W%C3%B6rterbuch:_Silbentrennung}}
erreichen.


%%%%%%%%%%%%%%%%%%%%%%%%%%%%%%%%%%%%%%%%%%%%%%%%%%%%%%%%%%%%
\subsection[Abkürzungen und Fachbegriff-Definitionen]{Abkürzungen und Fachbegriff-Definitionen}%
\index{Abkürzungen}%
\index{Akronyme}%
\index{Glossar}%
\index{Symbole}%
\label{sec:Abkuerzungen}
%%%%%%%%%%%%%%%%%%%%%%%%%%%%%%%%%%%%%%%%%%%%%%%%%%%%%%%%%%%%
%
In den Dateien
\printfilepathquoted{./preambel/Acronyms.tex},
\printfilepathquoted{./preambel/Glossary.tex} und
\printfilepathquoted{./preambel/GlossarySymbols.tex}
kann eine Liste der Abkürzungen, Fachbegriff"=Definitionen und Symbole angelegt werden.
Details hierzu finden sich im \cref{sec:Glossare}.


%%%%%%%%%%%%%%%%%%%%%%%%%%%%%%%%%%%%%%%%%%%%%%%%%%%%%%%%%%%%
\subsection[Frontmatter und Backmatter]{Frontmatter und Backmatter}%
\index{Frontmatter}%
\index{Backmatter}%
\index{Notation}%
\index{Symbolverzeichnis}%
\index{Literaturverzeichnis}%
\index{Abkürzungsverzeichnis}%
\index{Abbildungsverzeichnis}%
\index{Tabellenverzeichnis}%
\index{Listingsverzeichnis}%
\index{Glossar}%
\index{Stichwortverzeichnis}%
\label{sec:FrontmatterBackmatter}
%%%%%%%%%%%%%%%%%%%%%%%%%%%%%%%%%%%%%%%%%%%%%%%%%%%%%%%%%%%%
%
Die Einbindung von der Titelseite, Danksagung, dem englischen und deutschen Abstract,
einer Notationsübersicht und ggf. einem Symbolverzeichnis
erfolgt in der Datei \printfilepathquoted{./content/Inhalt-FrontMatter.tex}.
Die jeweilige Inhalte sind im Verzeichnis \printfilepathquoted{./content/00-FrontMatter/} zu finden.
Die Einbindung von Literatur-, Abbildungs-, Tabellen-, Quellcodelistings-, Abkürzungsverzeichnissen sowie ggf. einem Glossar, Stichwortverzeichnis (Index) und den Anhängen erfolgt in der Datei \printfilepathquoted{./content/Inhalt-BackMatter.tex}.

%%%%%%%%%%%%%%%%%%%%%%%%%%%%%%%%%%%%%%%%%%%%%%%%%%%%%%%%%%%%
\section[Grundsätzliches]{Grundsätzliches}%
\label{sec:Grundsätzliches}
%%%%%%%%%%%%%%%%%%%%%%%%%%%%%%%%%%%%%%%%%%%%%%%%%%%%%%%%%%%%
%
Beim Erstellen neuer Dateien \bzw Öffnen und Speichern bereits vorhandener Dateien ist darauf zu achten,
dass stets \index{UTF8-Kodierung}UTF-8 als Zeichenkodierung verwendet wird.
Wichtig ist dabei, dass alle tex-Dateien die UTF8-Kodierung ohne BOM haben,
worauf beim Anlegen neuer TeX-Dateien besonders zu achten ist
(am besten man kopiert und bearbeitet eine bereits vorhandene Datei).

Dank geeigneter Einstellungen in den Header-Dateien können deutsche Umlaute wie 
\index{Umlaute}%
ä,ö,ü,ß, Zeichen mit Akzent wie é sowie weitere UTF8-Zeichen wie \zB
\index{Anführungszeichen}%
\index{Sprache!unterschiedliche Anführungszeichen}%
„deutsche“, “englische”, »französische« oder «russische» Anführungszeichen 
direkt im Quellcode eingegeben werden ohne irgendwelche Umwege wie \zB 
\verb+"a+ für ä,
\verb+"u+ für ü,
\verb+\ss+, für ß und
\verb+'e+ für é.
Dies gilt insbesondere auch für Quellen des Literaturverzeichnisses (Bib-Dateien).

Die Zeiten, in welchen man sich bei der Eingabe deutscher Buchstaben und Anführungszeichen verkünsteln musste,
sind zum Glück endgültig vorbei.

\myexcl{Wichtig!}
Normale, gerade Anführungszeichen (\texttt{{\dq}}) haben in \LaTeX{} eine Sonderfunktion und sollten im Quellcode (außer in Listings) nicht verwendet werden.
Für die Eingabe von Anführungszeichen sollte man am besten die entsprechende Textstelle mit dem \lc{enquote\{...\}}"=Befehl umschließen.
Damit werden je nach Spracheinstellung des Dokumentes automatisch die richtigen Anführungszeichen gesetzt.
Außerdem werden so auch die \enquote{verschachtelten \enquote{Anführungszeichen}} korrekt behandelt.

Es empfiehlt sich, die einzelnen Sätze jeweils in einer neuen Zeile anzufangen.
Ein einfacher Zeilenumbruch wird von LaTeX wie ein Leerzeichen gehandhabt
und hat somit keinen Einfluss auf die Zeilenumbrüche im Ergebnisdokument.
Beim Rückwärtsspringen aus der PDF-Datei zum Quellcode wird dadurch jedoch
eine wesentlich bessere Lokalisierung der betroffenen Textstelle ermöglicht.
%%%%%%%%%%%%%%%%%%%%%%%%%%%%%%%%%%%%%%%%%%%%%%%%%%%%%%%%%%%%
\section[Globale Sprachumstellung und temporäre Sprachumschaltung]{Globale Sprachumstellung und \mbox{temporäre Sprachumschaltung} bei fremdsprachlichen Begriffen \mbox{und Textabschnitten}}%
\index{Fremdsprachen}%
\index{Sprache!Fremdsprache}%
\index{Sprache!globale Umstellung}%
\index{Sprache!temporäre Umschaltung}%
\label{sec:Sprache}
%%%%%%%%%%%%%%%%%%%%%%%%%%%%%%%%%%%%%%%%%%%%%%%%%%%%%%%%%%%%
%
Die Hauptsprache der Arbeit wird in der Datei \printfilepathquoted{./preambel/AlleSchalter.tex} festgelegt.
Aktuell werden nur Deutsch und Englisch als Hauptsprachen unterstützt.
Die Auswahl geschieht mit der Angabe des Wertes \printkeyword{true} oder \printkeyword{false}
in der Zeile \lstinline|\setUserDefinedBoolean{englishAsMainLanguage}{<Wert>}|.

Bei Verwendung von fremdsprachlichen Begriffen oder Textabschnitten
(\zB bei englischen oder französischen Zitaten in einer deutschsprachigen Arbeit
oder bei deutschen Begriffen in einer englischsprachigen Arbeit),
sollte man dies entsprechend markieren,
damit \LaTeX{} die richtigen Regeln für die \index{Silbentrennung}Silbentrennung
und die passenden Anführungszeichen bei Verwendung des Befehls
\lstinline|\enquote{...}| ansetzt.
Für die einzelnen Begriffe und kürzere Texte gibt es den Befehl
\lstinline|\foreignlanguage{Sprache}{...}|.
Dann wird für den Text in den geschweiften Klammern die in den eckigen Klammern angegebene Sprache verwendet.
Um die Sprache bis zum nächsten Aufruf des gleichen Kommandos dauerhaft umzustellen,
gibt es den Befehl \lstinline|\selectlanguage{Sprache}|.
Es gilt eine Liste der Sprachen aus dem Paket \pkg{babel}.
Für Deutsch sollte \printkeyword{ngerman} verwendet werden, was für die
\index{Rechtschreibung!neue deutsche}neue deutsche Rechtschreibung steht.

Für eine korrekte Behandlung der deutschen Kurzfassung und des englischen Abstracts
unabhängig von der gewählten Hauptsprache ist durch die Verwendung der Befehle
\lstinline|\textInGerman{...}| und \lstinline|\textInEnglish{...}| bereits gesorgt.
%%%%%%%%%%%%%%%%%%%%%%%%%%%%%%%%%%%%%%%%%%%%%%%%%%%%%%%%%%%%
\section[% Kurzversion für das Inhaltsverz., Kolumnentitel und PDF-Lesezeichen:
         Wichtiges zu Umbrüchen bei Überschriften
         \mbox{(Kurzversion für das Inhaltsverzeichnis \etc)}%
				]{% Langversion, die im Text gedruckt wird:
         Wichtiges zu Umbrüchen bei Überschriften
         (und ein Beispiel \newline für eine lange Überschrift,
         welche \newline für das Inhaltsverzeichnis und 
         \newline die Kolumnentiteln zu lang ist)%
				}%
\index{Zeilenumbruch}\index{Titel}\index{Ueberschrift@Überschrift}\index{PDF!Lesezeichen}%
\label{sec:Titles}
%%%%%%%%%%%%%%%%%%%%%%%%%%%%%%%%%%%%%%%%%%%%%%%%%%%%%%%%%%%%
%
%
Bei den Kapitelüberschriften kann man zwei Versionen definieren:
eine lange Überschrift in geschweiften Klammern, welche in der Arbeit selbst angezeigt wird, 
und optional eine Kurzversion in eckigen Klammern, welche im Inhaltsverzeichnis und in den 
\href{https://de.wikipedia.org/wiki/Kolumnentitel}{Kolumnentiteln}%
\footnote{Kolumnentitel sind Überschriften der einzelnen Seiten. Meist stehen sie in der Kopfzeile.}
angezeigt wird:

\begin{inlinelatex}
\section[Kurzversion]{Langer Titel}
\end{inlinelatex}

Dasselbe gilt für Bild- und Tabellenunterschriften. Hier kann man dem
\lc{caption}-Befehl ebenfalls einen optionalen Parameter in eckigen Klammern übergeben.

Um bei längeren Überschriften ein besseres Zeilenumbruchsverhalten im Inhaltsverzeichnis zu erzwingen
(\zB um eine ähnliche Länge aller Zeilen bei mehrzeiligen Überschriften zu erreichen),
ist manchmal eine Einfügung manueller Zeilenumbrüche in der Kurzversion der Überschriften erforderlich.
Allerdings funktioniert ein Einfügen manueller Zeilenumbrüche in der Kurzversion des Titels
etwa mit \verb+\newline+  leider nicht, da diese dann nicht nur im Inhaltsverzeichnis,
sondern auch in den Kolumnentiteln und PDF"=Lesezeichen zur Geltung kommen, 
was normalerweise nicht erwünscht ist.

Abhilfe schafft der folgende Trick:
man schließt den letzten, umzubrechenden Teil der Kurzversion des Titels in eine \verb+\mbox{}+.
Der Text, der in eine \verb+\mbox{}+ eingeschlossen wird, darf nicht umbrochen werden.
Im Kolumnentiteln und in den PDF"=Lesezeichen hat dies keine besondere Wirkung;
im Inhaltsverzeichnis führt dies jedoch dazu, dass \LaTeX{} den Zeilenumbruch vor der \verb+\mbox{}+ einfügt.
Ein entsprechendes Beispiel stellt die Überschrift dieses Abschnitts dar.
Dasselbe gilt für die ungünstig umbrochenen Wörter
(so empfiehlt der KSP Verlag darauf zu achten, dass keine seitenübergreifenden Wort-Trennungen und mehr als 4 nacheinanderfolgenden Zeilen mit Worttrennung am Ende der Zeile entstehen).
Weiteren Spielraum gibt das umbruchgeschützte Leerzeichen (eingefügt mit einer Tilde),
der Parameter \printkeyword{emergencystretch}, welcher global die maximal mögliche Dehnung des Zwischenraums zwischen den einzelnen Wörtern steuert,
sowie zusätzliche Steuermöglichkeiten der Silbertrennung, die auf Wikibooks%
\footnote{\url{https://de.wikibooks.org/wiki/LaTeX-W%C3%B6rterbuch:_Silbentrennung}}
beschrieben sind.
%%%%%%%%%%%%%%%%%%%%%%%%%%%%%%%%%%%%%%%%%%%%%%%%%%%%%%%%%%%%
\section{Aufzählungen}%
\index{Aufzählungen}%
\label{sec:Aufzählungen}
%%%%%%%%%%%%%%%%%%%%%%%%%%%%%%%%%%%%%%%%%%%%%%%%%%%%%%%%%%%%
%
Nachfolgend gibt es Beispiele für unterschiedliche Aufzählungen.
%
\begin{enumerate}
\item \justifying Aufzählungen mit längeren \enquote{Items}, die ganze Sätze enthalten, sollten mit \printkeyword{itemize} (unnummerierte Liste) \bzw \printkeyword{enumerate} (nummerierte Liste) eingefügt werden.

\item \justifying Dabei ist darauf zu achten, dass im Quellcode keine leeren Zeilen und damit kein zusätzlicher Absatz zwischen der Aufzählung und dem vorhergehenden Textabschnitt eingefügt wird.

\item \justifying Um die Aufzählung im Quellcode optisch vom vorangegangenem Textabschnitt zu trennen, können leere Zeilen mit einem Prozentzeichen verwendet werden, die von \LaTeX{} ignoriert werden.
\end{enumerate}
%
Bei Verwendung von \printkeyword{itemize} \bzw \printkeyword{enumerate} ist der vertikale Abstand zwischen den einzelnen \enquote{Items} ziemlich groß.
Dies ist gut für längere Punkte, sieht aber bei Aufzählungen mit kürzeren \enquote{Items} unschön aus.
Hierfür bieten die beiden Umgebungen \printkeyword{itemize*} und \printkeyword{enumerate*} mit einem Stern eine Abhilfe.
Das Ergebnis sind kompaktere Auflistungen mit kleinerem Abstand:
%
\begin{itemize*}
\item Punkt 1
\item Punkt 2
\item Punkt 3
\end{itemize*}

%%%%%%%%%%%%%%%%%%%%%%%%%%%%%%%%%%%%%%%%%%%%%%%%%%%%%%%%%%%%
\section[Bilder, Grafiken und Diagramme]{Bilder, Grafiken und Diagramme}
\label{sec:Bilder}
%%%%%%%%%%%%%%%%%%%%%%%%%%%%%%%%%%%%%%%%%%%%%%%%%%%%%%%%%%%%
%
Bei Einbindung von Grafiken sind zwei Fälle zu unterscheiden:
\begin{itemize*}
  \item reguläre \index{Bild!Binär-}Bilder in einem Binärformat
	      (\texttt{PNG}, \texttt{TIFF}, \texttt{JPG}, \texttt{PDF}, \etc)
	\item \index{Bild!Vektor-}Grafiken, die im \gls{gls:tikz}-Quellcode vorliegen 
\end{itemize*}

Grundlegender Unterschied bei der Einbindung \enquote{regulärer} Bilder
und \gls{gls:tikz}-Bilder ist, dass Binärformatgrafiken mit
\lc{includegraphics\{...\}} eingebunden werden,
während \gls{gls:tikz}-Grafiken mit \lc{input\{...\}}
eingebunden und von \printkeyword{latex} mitkompiliert werden.


%%%%%%%%%%%%%%%%%%%%%%%%%%%%%%%%%%%%%%%%%%%%%%%%%%%%%%%%%%%%
\subsection[Floats]{\index{Float}\index{Bild!Float}Floats}%
\label{sec:Floats}
%%%%%%%%%%%%%%%%%%%%%%%%%%%%%%%%%%%%%%%%%%%%%%%%%%%%%%%%%%%%
%
Üblicherweise werden Bilder, Tabellen und Codelistings in Fließumgebungen (floats) gesetzt,
damit \LaTeX\ sie frei im Text herumbewegen und dadurch geschickt positionieren kann.
Bei Bildern heißt die entsprechende Float-Umgebung \printkeyword{figure}, 
bei Tabellen \printkeyword{table}, bei Codelistings \printkeyword{lstlisting}.
Die Positionierung kann durch Angabe von Buchstaben
\printkeyword{h}, \printkeyword{t}, \printkeyword{b} und \printkeyword{p}
beeinflusst werden
(\enquote{\textbf{h}ere}, \enquote{\textbf{t}op}, \enquote{\textbf{b}ottom}, \enquote{extra \textbf{p}age}).
Eine genaue Auswirkung der Parameter
\printkeyword{h}, \printkeyword{t}, \printkeyword{b} und \printkeyword{p}
\bzw ihrer Kombinationen auf die Bildplatzierung ist jedoch nicht intuitiv.
Um diese zu verstehen, empfiehlt sich die Lektüre der Beschreibung
von Frank Mittelbach
\href{https://tex.stackexchange.com/questions/39017/how-to-influence-the-position-of-float-environments-like-figure-and-table-in-lat/39020#39020}{auf stackexchange.com}.%
\footnote{\url{https://tex.stackexchange.com/questions/39017/how-to-influence-the-position-of-float-environments-like-figure-and-table-in-lat/39020#39020}}

So versucht \LaTeX{} bei angegebener \printkeyword{t}- oder \printkeyword{b}-Option,
das Bild \bzw die Tabelle am oberen oder unteren Rand der Seite zu platzieren,
wodurch es passieren kann, dass der Absatz durch die Abbildung \etc
mitten in einem Satz aufgetrennt wird, was im professionellen Buchdruck unerwünscht ist.
Auch bei Einbindung von Floats
%mitten in einem Textabsatz
mit Hilfe der \printkeyword{h}-Option
platziert \LaTeX{} diese nicht unbedingt direkt an der Stelle der Einbindung im Quellcode.
%sondern zwischen zwei Zeilen in der Nähe der Einbindungsstelle, wobei auf Einhaltung
%möglichst gleichmäßiger Wortabstände im Absatz geachtet wird.
%Somit kann es dabei ebenfalls zu einer unerwünschten Auftrennung des Textes
%mitten in einem Satz kommen.
%
%Auch bei der Platzierung von Bildern und Tabellen mit Hilfe der \printkeyword{h}-Option
%in einem eigenen Absatz zwischen zwei Textabsätzen kann es zum unerwünschten Verhalten kommen.
%Wenn der verbleibende Platz die Platzierung eines Floats an der erwünschten Stelle nicht erlaubt,
%kann es schnell passieren, dass das betroffene Float auf die nächste Seite verschoben werden muss.
%Dabei wandelt \LaTeX{} die \printkeyword{h}-Option oft eigenwillig in die \printkeyword{ht}-Option um
%und verschiebt das Bild an den oberen Rand der Seite,
%%(beim Kompilieren des Quellcodes erhält man hierzu Feedback in Form einer Warnung), 
%wobei die oben beschriebene unerwünschte Auftrennung des Textes mitten in einem Satz entstehen kann.
%Muss ein Float soweit verschoben werden, dass das nächste Float von \LaTeX{} mitangefasst werden müsste, 
%wird das betroffene Float und alle Folgefloats unerwünschterweise ans Ende des Kapitels verschoben
%(genauer gesagt, an die Stelle, wo die nächste \lstinline|\clearpage|-Anweisung kommt).

Eine Abhilfe schafft eine konsequente Einbindung von Bildern und Tabellen mit Hilfe der \printkeyword{H}-Option,
%deren Nutzung in dieser Vorlage durch die Einbindung des \pkg{float}-Pakets aktiviert wurde und
welche eine Platzierung von Floats an der spezifizierten Stelle erzwingt.
Die Einbindung sollte als ein eigener Absatz zwischen zwei Textabsätzen erfolgen.

\myexcl{Wichtig!}
Wenn der verbleibende Platz die Platzierung eines Floats an der erwünschten Stelle nicht erlaubt,
wird dieses auf die nächste Seite verschoben. Dasselbe gilt für den nachfolgenden Inhalt.
Als Nebenwirkung können dabei auf der vorhergehenden Seite riesige leere Flächen entstehen.
Falls solche Nebeneffekte auftreten, müssen die betroffenen Bilder- und Tabellendefinitionen
manuell im Quellcode solange zwischen den Absätzen herumgeschoben werden,
bis eine optisch annehmbare Platzierung erreicht wird.
Dafür empfiehlt es sich, die einzelnen Bild- und Tabellendefinitionen in Extra-Dateien auszulagern,
so dass nur noch jeweils eine einzige Zeile mit der \lc{input\{...\}}-Anweisung herumgeschoben werden muss.
Prinzipiell empfiehlt sich eine endgültige Platzierung der Bilder und Tabellen erst ganz am Schluss,
nachdem alle anderen Korrekturen durchgeführt worden sind.

Wichtig anzumerken ist, dass alle Zeilen innerhalb der \printkeyword{figure}"=
\bzw der \printkeyword{table}"=Umgebung
mit einem Prozentzeichen abzuschließen sind.
Ansonsten werden überflüssige Leerzeichen eingefügt,
was zu unerwünschten Nebenwirkungen führen kann.

%%%%%%%%%%%%%%%%%%%%%%%%%%%%%%%%%%%%%%%%%%%%%%%%%%%%%%%%%%%%
\subsection[Binärbilder]{\index{Bild!Binär-}Binärbilder}%
\label{sec:Binaerbilder}
%%%%%%%%%%%%%%%%%%%%%%%%%%%%%%%%%%%%%%%%%%%%%%%%%%%%%%%%%%%%
%
Ein Beispiel für die Einbindung eines Bildes im Binärformat zeigt \cref{lst:binary-image}.
%Ein Beispiel für die Einbindung eines Bildes im Binärformat ist in \cref{lst:binary-image} angeführt.
%Man beachte die abschließenden Prozent-Zeichen am Ende jeder Zeile!

\begin{latex}[caption={Einbindung einer Binärgrafik in LaTeX},label={lst:binary-image}]
\begin{figure}[H]%
  \centering%
  \includegraphics[width=\linewidth]{Bildpfad/Dateiname}%
  \caption[Kurzversion für das Abbildungsverzeichnis]{%
           Eine tolle sehr lange Abbildungsunterschrift}%
  \label{fig:my-binary-image}%
\end{figure}
\end{latex}

Die Angabe des Pfades kann sowohl absolut als auch relativ
zum Verzeichnis der Hauptdatei oder zu einem der Pfade angegeben werden,
die in der Datei \printfilepathquoted{./preambel/AllePfade.tex} definiert sind.
Diese Pfade werden in angegebener Reihenfolge durchsucht.
Dasselbe gilt für die Dateierweiterung.
Ist keine Erweiterung definiert und liegen mehrere Bilder mit gleichem Namen jedoch unterschiedlicher Dateierweiterung vor,
wird die Reihenfolge, die in der Datei \printfilepathquoted{AllePfade.tex} definiert ist, verwendet.

Der \gls{ac:KSP} empfiehlt, Bilder und Grafiken auf die Satzspiegelbreite zu skalieren.
Dies wird durch die Option \printkeyword{width=\bs linewidth} umgesetzt.

Das Ergebnis sieht man in \cref{fig:my-binary-image}.

\begin{figure}[H]%
  \centering%
  \includegraphics[width=\linewidth]{./images/examples/img-example-01.jpg}%
  \caption[Einbindung eines Bildes]{Beispiel für die Einbindung eines Binärbildes}%
  \label{fig:my-binary-image}%
\end{figure}

Mit Hilfe von Paket \pkg{subfig} \cite{Cochran2005} können Bilder auch in
\index{Abbildung|see{Bild}}\index{Bild!Unterabbildung}Unterabbildungen gesetzt
und sowohl als ganzes (vgl. \cref{fig:subfloat-example}) als auch einzeln (vgl.
\cref{fig:subfloat-example-01,fig:subfloat-example-02,fig:subfloat-example-03,fig:subfloat-example-04})
referenziert werden.

\begin{figure}[H]%
	\centering%
	\subfloat[La Savoureuse, Lepuix, Frankreich][La Savoureuse, Lepuix, Frankreich \linebreak \textcopyright{} Thomas Bresson]{%
		\label{fig:subfloat-example-01}%
		\includegraphics[width=0.49\linewidth]{./images/examples/subfloat-example-01.jpg}%
	}%
	\hfill%
	\subfloat[Bangkok, Thailand][Bangkok, Thailand \linebreak \textcopyright{} Prachanart Viriyaraks]{%
		\label{fig:subfloat-example-02}%
		\includegraphics[width=0.49\linewidth]{./images/examples/subfloat-example-02.jpg}%
	}%
	\\\vspace{3mm}%
	\subfloat[Wahkeena Falls, Lincoln Park, USA][Wahkeena Falls, Lincoln Park, USA \linebreak \textcopyright{} srslyguys]{%
		\label{fig:subfloat-example-03}%
		\includegraphics[width=0.49\linewidth]{./images/examples/subfloat-example-03.jpg}%
	}%
	\hfill%
	\subfloat[Nacionalni park Plitvička jezer, Kroatien][Nacionalni park Plitvička jezer, Kroatien \linebreak \textcopyright{} Roman Bonnefoy]{%
		\label{fig:subfloat-example-04}%
		\includegraphics[width=0.49\linewidth]{./images/examples/subfloat-example-04.jpg}%
	}%
	\caption[Bild mit Unterabbildungen]{Wasserfälle der Welt als Beispiel für Unterabbildungen}%
	\label{fig:subfloat-example}%
\end{figure}

Der Beispielcode dafür ist in \cref{lst:subfigures} dargestellt.

\begin{latex}[caption={Unterabbildungen in LaTeX},label={lst:subfigures}]
\begin{figure}[H]%
	\centering%
	\subfloat[Kurzversion][Unterbezeichnung 1]{%
		\label{fig:UnterAbb1}%
		\includegraphics[width=0.49\linewidth]{Bildpfad/Bild1}%
	}%
	\hfill%
	\subfloat[Kurzversion][Unterbezeichnung 2]{%
		\label{fig:UnterAbb2}%
		\includegraphics[width=0.49\linewidth]{Bildpfad/Bild2}%
	}%
	\\\vspace{3mm}%
	\subfloat[Kurzversion][Unterbezeichnung 3]{%
		\label{fig:UnterAbb3}%
		\includegraphics[width=0.49\linewidth]{Bildpfad/Bild3}%
	}%
	\hfill%
	\subfloat[Kurzversion][Unterbezeichnung 4]{%
		\label{UnterAbb4}%
		\includegraphics[width=0.49\linewidth]{Bildpfad/Bild4}%
	}%
\caption[Kurzversion]{Langversion der Bildunterschrift}%
\label{fig:MeinGanzesBild}%
\end{figure}
\end{latex}


%%%%%%%%%%%%%%%%%%%%%%%%%%%%%%%%%%%%%%%%%%%%%%%%%%%%%%%%%%%%
\subsection[TikZ-Grafiken]{\index{TikZ}\index{Bild!TikZ}\gls{gls:tikz}-Grafiken}%
\label{sec:TikZ}
%%%%%%%%%%%%%%%%%%%%%%%%%%%%%%%%%%%%%%%%%%%%%%%%%%%%%%%%%%%%
%
\Gls{gls:tikz} eignet sich hervorragend, um wissenschaftliche Zeichnungen,
Vektorgrafiken und \index{Diagramm}Diagramme direkt mithilfe von LaTeX
zu setzen, sodass die Schrift direkt zum restlichen Dokument passt.
Zu dem \pkg{tikz}-\gls{gls:pkg} und dem darauf aufsetzenden \pkg{PGFplots}-\gls{gls:pkg} gibt
es hervorragende Dokumentationen \parencites{Tantau2013}{Feuersaenger2014}.
Mit \gls{gls:tikz} und \gls{gls:pgfplots} lassen sich viele gute Sachen machen.
%
%\setcounter{figure}{10}

Der Code für die Einbindung einer \gls{gls:tikz}-Grafik steht in \cref{lst:tikz-figure}.
\begin{latex}[caption={Einbindung einer TikZ-Zeichnung in LaTeX},label={lst:tikz-figure}]
\begin{figure}[H]%
  \centering%
  \tikzsetnextfilename{TikZ-Bild}%
  \resizebox{\textwidth}{!}{%   <--- optionale Skalierung
    \input{./figures-src/TikZ-Bild.tex}%
  }%                            <--- optionale Skalierung
  \caption[Kurzversion für das Abbildungsverzeichnis]{%
           Eine tolle sehr lange Abbildungsunterschrift}%
  \label{fig:my-tikz-figure}%
\end{figure}
\end{latex}

Eine Skalierung auf die volle Seitenbreite oder ein vielfaches davon im Falle von Unterabbildungen kann bei Bedarf mit Hilfe der Anweisung
\lstinline|\resizebox{\textwidth}{!}{...}|
durchgeführt werden.

Das Kommando \lstinline|\tikzsetnextfilename{...}| ist nicht unbedingt notwendig,
 aber sehr zu empfehlen, da dies als Name für das temporäre Kompilat im Ordner
\printfilepathquoted{./figures-compiled/} genommen wird.
Dieser sollte gleich dem Namen der Quelldatei (ohne Endung) gewählt werden.
Ansonsten nimmt \texttt{pdflatex} eine hochlaufende Nummer als Dateiname,
was die Fehlersuche sehr erschwert.

Nachfolgend finden sich einige Beispiele für TikZ-Zeichnungen, nämlich
ein kommutatives Diagramm (\cref{fig:kpca}),
eine Übersicht über die KIT-Corporate-Identity-Farben (\cref{fig:kit-colors}),
ein Netzwerkkommunikationsgraph (\cref{fig:net-comm}),
einfache \index{Diagramm!Punkt-}Punktdiagramme (\cref{fig:ica})
und etwas aufwendigere Diagramme mit mehreren
\index{Achsensystem|see{Diagramm}}Achsensystemen (\cref{fig:pca}).

Vorzugsweise sollten für die Grafiken die KIT-Corporate-Identity-Farben verwendet werden,
die sowohl eine RGB"=Definition für die Darstellung online,
als auch eine CMYK"=Definition für den Offset-Druck haben.%
\footnote{Bei Erstellung der Manuskriptversionen
für die Online-Veröffentlichung \bzw für den Offset-Druck
ist auf die korrekte Einstellung der Option \printkeyword{useCMYKcolors}
in der Datei \printfilepathquoted{preambel/AlleSchalter.tex} zu achten
(\printkeyword{false} für die Online-Veröffentlichung und \printkeyword{true} für den Druck!}

\begin{figure}[H]%
	\centering%
  \tikzsetnextfilename{kpca}%
	%\resizebox{\textwidth}{!}{%
		\input{./figures-src/kpca.tex}%
	%}%
	\caption{Kommutative Diagramm mit TikZ}%
  \label{fig:kpca}%
\end{figure}

\begin{figure}[H]%
	\centering%
  \tikzsetnextfilename{kit-colors}%
	%\resizebox{\textwidth}{!}{%
		\input{./figures-src/kit-colors.tex}%
	%}%
	\caption{KIT-Corporate-Identity-Farben}%
  \label{fig:kit-colors}%
\end{figure}

\begin{figure}[H]%
	\centering%
  \tikzsetnextfilename{net-comm}%
	\resizebox{0.75\textwidth}{!}{%
		\input{./figures-src/net-comm.tex}%
	}%
	\caption{Netzwerkkommunikationsgraph mit TikZ}%
  \label{fig:net-comm}%
\end{figure}

\begin{figure}[H]%
	\centering%
	\subfloat[Ursprünglicher Merkmalsraum]{%
		\label{fig:ica-1}%
		\tikzsetnextfilename{ica-1}%
		\resizebox{0.45\textwidth}{!}{%
			\input{./figures-src/ica-1.tex}
		}%
	}%
	\hfill%
	\subfloat[Transienter Merkmalsraum (Nach Whitening, z.\,B. durch \gls{ac:PCA} inkl. Normalisierung)]{%
		\label{fig:ica-2}%
		\tikzsetnextfilename{ica-2}%
		\resizebox{0.45\textwidth}{!}{%
			\input{./figures-src/ica-2.tex}%
		}%
	}%
	\hfill%
	\subfloat[Transformierter Merkmalsraum]{%
		\label{fig:ica-3}%
		\tikzsetnextfilename{ica-3}%
		\resizebox{0.45\textwidth}{!}{%
			\input{./figures-src/ica-3.tex}%
		}%
}%
\caption{Diagramme mit TikZ direkt in LaTeX (hier: Die Schritte der \enquote{Independent component analysis})}%
\label{fig:ica}%
\end{figure}

\begin{figure}[H]
	\centering%
	\subfloat[Ungünstige Projektion]{%
		\label{fig:pca-1}%
		\tikzsetnextfilename{pca-1}%
		\resizebox{0.48\textwidth}{!}{%
			\input{./figures-src/pca-1.tex}%
		}%
	}\hfill%
	\subfloat[Zielführende Projektion]{%
		\label{fig:pca-2}%
		\tikzsetnextfilename{pca-2}%
		\resizebox{0.48\textwidth}{!}{%
			\input{./figures-src/pca-2.tex}%
		}%
	}%
	\caption{Aufwändiges Diagramm mit TikZ (hier: Probleme der \enquote{Principal component analysis})}%
	\label{fig:pca}%
\end{figure}
%%%%%%%%%%%%%%%%%%%%%%%%%%%%%%%%%%%%%%%%%%%%%%%%%%%%%%%%%%%%
\section{Tabellen}%
\label{sec:Tabellen}
%%%%%%%%%%%%%%%%%%%%%%%%%%%%%%%%%%%%%%%%%%%%%%%%%%%%%%%%%%%%
Für eine ausführliche Erläuterung auch über gute und schlechte Tabellen
und deren Gestaltung empfiehlt sich die Lektüre des entsprechenden Artikels im \LaTeX{}"=Kompendium auf Wikibooks%
\footnote{\url{https://de.wikibooks.org/wiki/LaTeX-Kompendium:_Tabellen}}
sowie die Dokumentation des \pkg{booktabs}-Pakets \cite{Fear2005}.
\Ua stellt dieses Paket die Befehle
\begin{itemize*}
	\item \lstinline|\toprule|
	\item \lstinline|\midrule|
	\item \lstinline|\bottomrule|
\end{itemize*}
zur Verfügung.

Der Code für eine einfache Tabelle ist im \cref{lst:tabellenbeispiel} angeführt.
Damit alle Tabellen einheitlich aussehen, empfiehlt es sich,
überall das gleiche Makro zur Parameteranpassung zu verwenden.
Dazu könnte beispielsweise das vordefinierte Makro
\lstinline|\tablestylecommon| verwendet werden, welches in der Datei
\printfilepathquoted{./preambel/TableCommands.tex} definiert ist.
Zu beachten ist, dass die Anpassungen an den Tabellenparametern wie \zB
\begin{itemize*}
\item die Verkleinerung der Schriftgröße,
\item die Entfernung der seitlichen Ränder sowie
\item die Erhöhung der Abstandes zwischen den einzelnen Tabellenzeilen
\end{itemize*}
innerhalb einer Umgebung gesetzt werden müssen, damit sie lokal wirken und
sich nicht auf die nachfolgenden Tabellen erstrecken.
Wird die Tabelle nicht innerhalb einer \printkeyword{table}-Umgebung
\oae gesetzt, muss man eine eigene Umgebung mit
\lstinline|\begingroup| ... \lstinline|\endgroup|
definieren.

\begin{latex}[caption={Einfache Tabelle in \LaTeX},label={lst:tabellenbeispiel}]
\begin{table}%
	\tablestylecommon% Aussehen anpassen
	\caption{Dateien der Vorlage}%
	\label{tab:tabellenbeispiel}%
	\begin{tabular}{l l l}%
		\toprule%
		Datei       &  Bedeutung    &  Benutzerinteraktion \\%
		\midrule%
		Diss.tex  &  Hauptdatei   &  nein     \\%
		images/   &  Bilder       &  ja       \\%
		content/  &  Kapitel      &  ja       \\%
		\bottomrule%
	\end{tabular}%
\end{table}
\end{latex}

Das Ergebnis sieht man in \cref{tab:tabellenbeispiel}.

\begin{table}%
	\tablestylecommon% Aussehen anpassen
	\caption{Dateien der Vorlage}%
	\label{tab:tabellenbeispiel}%
	\begin{tabular}{l l l}%
		\toprule%
		Datei       &  Bedeutung    &  Benutzerinteraktion \\%
		\midrule%
		Diss.tex  &  Hauptdatei   &  nein     \\%
		images/   &  Bilder       &  ja       \\%
		content/  &  Kapitel      &  ja       \\%
		\bottomrule%
	\end{tabular}%
\end{table}

Ist die Tabelle seitenbreit, sollte statt der \printkeyword{tabular}-Umgebung die \printkeyword{tabularx}-Umgebung verwendet werden (\cref{tab:AlternierendeZeilenfarben}).
Mit dem \printkeyword{tablealtcolored}-Makro kann man die einzelnen Zeilen mit alternierender Farbe einfärben.
Die Anpassungsmöglichkeiten für das Aussehen der Tabellen befinden sich in der Datei \printfilepathquoted{./preambel/TableCommands.tex}.

%% Breite der ersten und zweiten Spalte auf Inhalt anpassen --> "l".
%% Breite der dritten Spalte automatisch bestimmen --> "X" (oder "x", "F", "y", "Y")
%% "X" = Blocksatz (mit Silbentrennung),
%% "x" = linksbündiger Flattersatz ohne Silbentrennung
%% "F" = linksbündiger Flattersatz mit Silbentrennung (KSP-Verlag: unerwünscht),
%% "y" = rechtsbündiger Flattersatz ohne Silbertrennung
%% "Y" = rechtsbündiger Flattersatz mit Silbertrennung
\begin{table}%
\tablestylealtcolored%
\caption{Seitenbreite Tabelle mit alternierender Zeilenfarbe}%
\label{tab:AlternierendeZeilenfarben}%
	\begin{tabularx}{\linewidth}{llX}%
		\toprule%
		\rowcolor{tableheadcolor}%
		\tableheadfont Tabellenkopf &	\tableheadfont Tabellenkopf & \tableheadfont Tabellenkopf 
		\tabularnewline%
		\midrule%
		Inhalt  & Etwas längerer Inhalt & Weiterer Inhalt \tabularnewline
		Inhalt  & Etwas längerer Inhalt & Weiterer Inhalt \tabularnewline
		Inhalt  & Etwas längerer Inhalt & Weiterer Inhalt \tabularnewline
		%% Zwischenkopf ---------------------------------------------
		\rowcolor{tablesubheadcolor}%
		\multicolumn{3}{l}{\bfseries\color{KITblue} Zwischenkopfzeile}%
		\tabularnewline%
		%%-----------------------------------------------------------
		Inhalt  & Etwas längerer Inhalt & Weiterer Inhalt \tabularnewline
		Inhalt  & Etwas längerer Inhalt & Weiterer Inhalt \tabularnewline
		Inhalt  & Etwas längerer Inhalt & Weiterer Inhalt \tabularnewline
		%% Zwischenkopf ---------------------------------------------
		\rowcolor{tablesubheadcolor}%
		\multicolumn{3}{l}{\bfseries\color{KITgreen} Zwischenkopfzeile}%
		\tabularnewline
		%%-----------------------------------------------------------
		Inhalt  & Etwas längerer Inhalt & Weiterer Inhalt \tabularnewline%
		Inhalt  & Etwas längerer Inhalt & Weiterer Inhalt \tabularnewline%
		\bottomrule%
	\end{tabularx}%
\end{table}

%% ------------------------------------------------------------
Lange Tabellen, die über mehrere Seiten gehen und umbrochen werden sollen, können mit
\lstinline|\LTXtable{\textwidth}{Datei}|
eingebunden werden, wobei die Tabelle in eine eigene Datei ausgelagert werden muss.
Ein Beispiel dafür sieht man in \cref{tab:MehrseitigeTabelle}.

% Angaben zur Tabellenstil innerhalb einer eigenen Umgebung, damit diese nur lokal wirken
\begingroup%
\tablestylealtcolored%
\setlength{\topskip}{0pt}%
\setlength{\LTpre}{0pt}%
\setlength{\LTpost}{0pt}%
\LTXtable{\textwidth}{tables/LongTableExample.tex}
\endgroup%

Sollte eine Tabelle so breit sein, dass sie nicht mehr horizontal auf
eine Seite passt, so ist es natürlich möglich, diese mithilfe des Pakets
\printkeyword{rotfloat} \parencite{Sommerfeldt2004} in eine
\printkeyword{sidewaystable} statt in eine \printkeyword{table}-Umgebung zu setzen.
Also so:
\begin{latex}[caption={Gedrehte Tabelle},label={lst:rotated-table}]
\begin{sidewaystable}
  \raggedleft% rechtsbündig, um die Tabelle oben auszurichten
  \begin{tabular}{...}%
    ...
  \end{tabular}%
  \caption{Bezeichnung}%
  \label{Referenzmarke}%
\end{sidewaystable}%
\end{latex}

Das Ergebnis sieht man in \cref{tab:ex-sideways}.

%% Um 90° gedrehte Tabelle
%
\begin{sidewaystable}[p]
\raggedleft%% Rechtsbündig, damit die Tabelle oben ausgerichtet wird (Forderung KSP)
\renewcommand{\arraystretch}{1.3}%% Etwas mehr Abstand zwischen den Tabellenzeilen
%% Tabelle mit einer festen Breite (hier: so breit wie die Seitenhöhe).
%% Spaltendefinition:
%% Erste Spalte vom Typ "l": Spalte so breit, wie der Inhalt, linksbündige Ausrichtung
%% Zweite Spalte vom Typ "L": Spalte mit einer festen Breitenangabe, linksbündiger Flattersatz. Inhalt wird umbrochen, wenn zu lang
%% 2 Spalten vom Typ "x":  die Breite wird automatisch berechnet
%% in Abhängigkeit von der Gesamtbreite der Tabelle und Breite anderer Spalten,
%% Inhalt wird umbrochen, wenn zu lang
%% Format: linksbündiger Flattersatz ohne Silbentrennung, 
%% Weitere Möglichkeiten:
%% "X" = Blocksatz (mit Silbentrennung)
%% "F" = linksbündiger Flattersatz mit Silbentrennung (KSP-Verlag: unerwünscht),
%% "y" = rechtsbündiger Flattersatz ohne Silbertrennung
%% "Y" = rechtsbündiger Flattersatz mit Silbertrennung
\begin{tabularx}{\linewidth}{l L{10em} *{2}x}%
	\toprule%
	\rowcolor{tableheadcolor}%
	\tableheadfont Spalte vom Typ \enquote{l} &	
	\tableheadfont Spalte vom Typ \enquote{L} & 
	\tableheadfont Spalte vom Typ \enquote{x} & 
	\tableheadfont Spalte vom Typ \enquote{x}
	\tabularnewline%
	\midrule%
		Spaltenbreite = Breite des Inhalts & 
		Spalte mit der festen Breite von 10\,em & 
		Spaltenbreite wird automatisch bestimmt. Inhalt wird bei Bedarf umbrochen. & 
		Spaltenbreite wird automatisch bestimmt. Inhalt wird bei Bedarf umbrochen.
	\tabularnewline
	%%-----------------------------------------------------------
	Nächste Zeile & Etwas längerer Inhalt & Weiterer Inhalt & und noch etwas \tabularnewline
	Weitere Zeile & Etwas längerer Inhalt & Weiterer Inhalt & und noch etwas \tabularnewline
	%% Zwischenkopf ---------------------------------------------
	\rowcolor{tablesubheadcolor}%
	\multicolumn{4}{l}{\bfseries\color{KITblue} Zwischenkopfzeile, welche über mehrere Spalten geht}%
	\tabularnewline%
	%%-----------------------------------------------------------
	Inhalt  & Etwas längerer Inhalt & Weiterer Inhalt & und noch etwas \tabularnewline
	Inhalt  & Etwas längerer Inhalt & Weiterer Inhalt & und noch etwas \tabularnewline
	Inhalt  & Etwas längerer Inhalt & Weiterer Inhalt & und noch etwas \tabularnewline
	%% Zwischenkopf ---------------------------------------------
	\rowcolor{tablesubheadcolor}%
	\multicolumn{4}{l}{\bfseries\color{KITgreen} Zwischenkopfzeile, welche über mehrere Spalten geht}%
	\tabularnewline
	%%-----------------------------------------------------------
	Inhalt  & Etwas längerer Inhalt & Weiterer Inhalt & und noch etwas \tabularnewline
	Inhalt  & Etwas längerer Inhalt & Weiterer Inhalt & und noch etwas \tabularnewline
	\bottomrule%
\end{tabularx}%
\caption{Beispiel für eine breite, gedrehte Tabelle}%
\label{tab:ex-sideways}%
\end{sidewaystable}
%%%%%%%%%%%%%%%%%%%%%%%%%%%%%%%%%%%%%%%%%%%%%%%%%%%%%%%%%%%%
\section{Mathematische Sätze, Lemmas, Definitionen etc.}%
\label{sec:Theoreme}
%%%%%%%%%%%%%%%%%%%%%%%%%%%%%%%%%%%%%%%%%%%%%%%%%%%%%%%%%%%%
%
Für eine mathematische Ausarbeitung gibt es LaTeX-\glspl{gls:umgebung}, um
\index{Satz|see{Theorem}}\index{Theorem}Sätze (Theoreme), \index{Lemma|see{Theorem}}Lemma,
\index{Beispiel|see{Theorem}}Beispiele \etc im üblichen Stil von
Mathematik-Büchern zu setzen und zu referenzieren. Vordefiniert sind die
\glspl{gls:umgebung}
\begin{itemize*}
  \item \texttt{theorem} für Sätze
  \item \texttt{definition} für Definitionen
  \item \texttt{lemma} für Lemma
  \item \texttt{corollary} für Korollare
  \item \texttt{proposition} für Propositionen
\end{itemize*}
Die übliche Verwendung ist
\begin{latex}[caption={Beispiel für Theorem-Umgebungen},label={lst:ntheorem}]
\begin{theorem}[Optionaler Name]\label{thm:my-theorem}
...
\end{theorem}
\end{latex}
Weitere Informationen findet man in der Dokumentation zum \texttt{ntheorem}-Paket
\parencite{May2011}. Das Ganze sieht dann beispielsweise wie folgt aus.

\begin{theorem}[Theorem von Arthur Dent]
\label{thm:arthur-dent} Die Antwort auf die Frage nach dem Leben, dem Universum und den ganzen Rest ist 42.
\end{theorem}

\begin{definition}
\LaTeX{} ist eine von Leslie Lamport 1980 entwickelter Satz von Makros zur Erweiterung von \TeX.
\end{definition}

\begin{proposition}[Zweifelhafte Folgerung]
LaTeX ist schön. Beweis folgt unmittelbar aus \cref{thm:arthur-dent}.
\end{proposition}

%%%%%%%%%%%%%%%%%%%%%%%%%%%%%%%%%%%%%%%%%%%%%%%%%%%%%%%%%%%%
\section{Quellcode-Listings}%
\index{Listing!Gestaltungsstil}%
\label{sec:Listings}
%%%%%%%%%%%%%%%%%%%%%%%%%%%%%%%%%%%%%%%%%%%%%%%%%%%%%%%%%%%%
%
Zum Einbinden und Formatieren von \index{Code|see{Listing}}Quellcode"=Beispielen
-- sog. \index{Listing}Listings -- wird das Paket \pkg{listings}
\parencite{Hoffmann2014} verwendet.
Das Hervorheben von \index{Schlusselwort@Schlüsselwort}Schlüsselwörtern
wird von LaTeX automatisch erledigt,
wenn die korrekte Sprache des Listings angegeben ist.
Dies geschieht mit Hilfe der Option \printkeyword{language}
oder durch die Angabe eines entsprechend für die jeweilige Sprache
vordefinierten Gestaltungsstils.

Im Befehl \lstinline|\lstset{...}|,
welcher in der Datei \printfilepathquoted{./preambel/preambel.tex} zu finden ist,
kann man einen globalen Stil für alle Listings vorgeben
(welcher jedoch bei Bedarf im Einzelfall überrufen werden kann).
Aktuell ist der etwas weiter oben im Code mit dem Befehl
\lstinline|\lstdefinestyle{...}| vordefinierte Stil
\printkeyword{numbersframed} als Standardgestaltungsstil ausgewählt.
Möchte man die Zeilennummern nicht anzeigen, wäre stattdessen
\printkeyword{nonumbersframed} zu wählen.

%Neben diesen beiden Gestaltungsstilen sind in der Datei
%\printfilepathquoted{./preambel/preambel.tex}
%auch noch \printkeyword{python}, \printkeyword{java}, \printkeyword{C++}
%und \printkeyword{latex} als Gestaltungsstile vordefiniert.
%Bei Bedarf lassen sich dort weitere Stile definieren und auswählen.

Zur Vereinfachung wurden für die \og Programmiersprachen
eigene Umgebungen definiert, die mit \lstinline|\begin{<name>}...\end{<name>}|
direkt verwendet werden können.
Ein Beispiel für die Einbindung eines Java"=Quellcodelistings mit Hilfe einer solchen Umgebung ist in \cref{lst:java-listing} gezeigt.

Man beachte, dass anders als bei Abbildungen und Tabellen
die Bezeichnung (\texttt{caption}) und die Referenzmarke (\texttt{label})
nicht als gesonderte Befehle, sondern als optionale Argumente übergeben werden.
Dies liegt daran, dass ein Listing in der Regel keine Fließumgebung ist,
sondern an der Stelle im Text erscheint, an der sie im Code auch steht.
Ferner folgt ein Listing den ganz normalen Seitenumbruchsregeln.
Das heißt, ein überlanger Code wird einfach umgebrochen. 
Um ein Listing zu einem Fließobjekt zu machen, muss das optionale Argument
\texttt{float=<tbp>} angegeben werden.
Die \index{Platzierung}Platzierungsangabe \enquote{\texttt{h}} für \enquote{hier}
ist nicht erlaubt. Denn dies ist das Standardverhalten ohne \texttt{float}.
%
\begin{latex}[caption={Beispiel eines LaTeX-Quellcodes für ein Java-Listing},label={lst:java-listing}]
\begin{java}[caption={Ein Hello-World Beispiel in Java},
             label={lst:hello-world}]
// Klasse HelloWorld
public class HelloWorld {
  public static void main( String[] args ) {
    System.out.println( "HelloWorld" );
  }
}
\end{java}
\end{latex}
%
Dieses Listing bewirkt das folgende Ergebnis:
%
\begin{java}[caption={Ein Hello-World Beispiel in Java},
             label={lst:hello-world}]
// Klasse HelloWorld
public class HelloWorld {
  public static void main( String[] args ) {
    System.out.println( "HelloWorld" );
  }
}
\end{java}


%%%%%%%%%%%%%%%%%%%%%%%%%%%%%%%%%%%%%%%%%%%%%%%%%%%%%%%%%%%%
\section{Mathematik}%
\label{sec:Mathe}
%%%%%%%%%%%%%%%%%%%%%%%%%%%%%%%%%%%%%%%%%%%%%%%%%%%%%%%%%%%%
%
Grundsätzlich gilt alles, was in \parencites{ams1999a}{ams1999b} steht. 
Um eine einheitliche Typografie von
\index{Skalare}Skalaren, \index{Vektoren}Vektoren, \index{Matrizen}Matrizen,
\index{Zufallsvariablen}Zufallsvariablen \etc zu vereinfachen, ist in der Datei
\printfilepathquoted{preambel/Math.tex} eine Menge Kurzkommandos definiert.
In diese Datei einfach mal reinschauen, welche Kurzkommandos es gibt.

Grundsätzlich empfiehlt es sich, eigene Mathe-Makros zu definieren und durchgehend zu verwenden.
Damit kann man bei Bedarf die Formatierung der Variablen an einer einzigen Stelle anpassen,
ohne den gesamten Code des Arbeit durchgehen zu müssen.

Die Einbettung der Formeln geschieht mit der \printkeyword{equation}"=Umgebung.
Da die Formeln normalerweise keine eigenständigen Absätze darstellen,
sollen zwischen dem Fließtext und der Formeldefinition keine leeren Zeilen stehen,
da \LaTeX{} diese als Absätze behandelt und unnötigerweise zusätzliche vertikale Abstände hinzufügt.

Auch sollte man darauf achten, dass zwischen einer Zahl und der physikalischen Einheit ein umbruchgeschütztes Leerzeichen verwendet wird.
Allerdings sollte statt dem üblichen (dehnbaren) umbruchgeschützen Wortzwischenraum, 
welches mit einer Tilde erzeugt wird, ein schmälerer Abstand verwendet werden.
Dieser wird mit \verb+\,+  erzeugt und beträgt standardmäßig 0,1667\,em. % Dies geht so:
%
%\begin{inlinelatex}
%Der eingefügte Abstand beträgt 0,1667\,em.
%\end{inlinelatex}

Auf zwei besondere Kommandos wird näher eingegangen,
weil dies in \LaTeX{} häufig falsch gemacht wird.
\begin{itemize}
  \item Für die Matrixtransponierte gibt es das Kommando \verb#\Tr#, also
	\verb#$A^{\Tr}$# liefert $A^{\Tr}$	
	\item \justifying Bei \index{Integral}Integralen muss das \enquote{Differential-d} gemäß
	ISO in aufrechter Schrift als Operator gesetzt sein mit einem kleinen Abstand
	zum Integranden. Hierfür gibt es das spezielle Kommando \verb#\diff#.
	\Dhe, die Schreibweise \verb#\int^1_0 x^2 d x = \frac{1}{3}# (\cref{eq:Diff-wrong}) ist falsch,
	während \verb#\int^1_0 x^2 \diff x = \frac{1}{3}# (\cref{eq:Diff-correct}) das Richtige liefert:
\end{itemize}

\begin{equation}
	\int^1_0 x^2 d x = \frac{1}{3} \qquad \text{(falsche Typografie!)}
	\label{eq:Diff-wrong}
\end{equation}

\begin{equation}
	\int^1_0 x^2 \diff x = \frac{1}{3} \qquad \text{(richtige Typografie!)}
	\label{eq:Diff-correct}
\end{equation}


%%%%%%%%%%%%%%%%%%%%%%%%%%%%%%%%%%%%%%%%%%%%%%%%%%%%%%%%%%%%
\section{Querverweise und Hyperlinks}%
\label{sec:Querverweise}
%%%%%%%%%%%%%%%%%%%%%%%%%%%%%%%%%%%%%%%%%%%%%%%%%%%%%%%%%%%%
%
Querverweise sollten nicht mit dem Befehl \lstinline|\ref{...}| gesetzt werden,
sondern mit \lstinline|\cref{...}| und verwandten Befehlen aus dem Paket
\pkg{cleveref} \cite{Cubitt2013}.
Diese Befehle haben den Vorteil nicht nur die Nummer zu referenzieren,
sondern auch den Typ mit anzugeben.
Hinzu kommt eine intelligente Verwendung der Pluralform und
\index{Sortierung}Sortierung bei Mehrfachaufzählungen auch unterschiedlichen Typs.
Will man \bspw auf zwei Abbildungen und eine Tabelle mit den Marken (\enquote{Labels})
%
\begin{itemize*}
\item \texttt{fig:subfloat-example}
\item \texttt{tab:files-dirs-of-template}
\item \texttt{fig:kit-colors}
\end{itemize*}
%
verweisen, so schreibt man einfach per Komma getrennt
%
\begin{latex}[caption={Cleveres Referenzieren mit \bs cref},label={lst:cref}]
\cref{fig:subfloat-example,
      tab:ex-sideways,
      fig:kit-colors}
\end{latex}
%
und erhält als Resultat
\enquote{\cref{fig:subfloat-example,tab:ex-sideways,fig:kit-colors}}.

\index{URL}\index{Internetadresse|see{URL}}Internetadressen werden in das Kommando \lstinline|\url{...}| eingefasst.
Außerdem besteht die Möglichkeit mit dem Befehl \lstinline|\href{<URL>}{Text}| eine Textstelle mit einem Hyperlink zu versehen.
%%%%%%%%%%%%%%%%%%%%%%%%%%%%%%%%%%%%%%%%%%%%%%%%%%%%%%%%%%%%
\section{Literaturverzeichnisse}%
\label{sec:Literaturverzeichnisse}%
\index{Bibliografie}\index{Literaturverzeichnis}
%%%%%%%%%%%%%%%%%%%%%%%%%%%%%%%%%%%%%%%%%%%%%%%%%%%%%%%%%%%%
%
Defaultmäßig befinden sich alle Literaturangaben in der Datei
\printfilepathquoted{bib/Diss.bib}.
Falls ein davon abweichender Pfad oder zusätzliche Literaturdatenbanken
verwendet werden sollen, können entsprechende Pfade in der Datei
\printfilepathquoted{preambel/AllePfade.tex} mit
\lstinline|\addbibresource{...}|
festgelegt werden.
Bitte beachten, dass bei den Literaturdatenbanken kein \gls{gls:bibtex}-,
sondern \gls{gls:biblatex}-Format verwendet werden soll,
also Umlaute direkt eingegeben werden können.

Die Vorlage ermöglicht es, neben dem Haupt"=Literaturverzeichnis
jeweils separat durchnummerierte Listen der eigenen Publikationen und Patente
sowie der betreuten studentischen Arbeiten zu erstellen.
Standardmäßig werden alle Literaturangaben aus der Hauptliteraturdatenbank
\printfilepathquoted{bib/Diss.bib} extrahiert.
Durch Änderung der Definitionen der Pfade
\printkeyword{bibpathOwnPatents},
\printkeyword{bibpathOwnPublications} und
\printkeyword{bibpathStudentTheses}
in der Datei
\printfilepathquoted{preambel/AllePfade.tex}
können jedoch für die einzelnen Auflistungen 
jeweils eigene, speziell dafür angelegte Datenbanken verwendet werden.

Die Einbindung der Literaturverzeichnisse in das Manuskript erfolgt in der Datei
\printfilepathquoted{content/Inhalt-BackMatter.tex}.
Die Zusatz"=Literaturverzeichnisse werden separat erstellt und sind vom Haupt"=Literaturverzeichnis unabhängig.
Dazu werden in den entsprechenden Abschnitten der Datei,
welche jeweils mit Hilfe der \printkeyword{refsection}"=Umgebung abgegrenzt sind,
die kompletten Inhalte der jeweiligen Literaturdatenbank unsichtbar referenziert.
Dies geschieht mit dem \lstinline|\nocite{*}|"=Befehl.
Alternativ dazu können die einzelnen Werke mit
\lstinline|\nocite{Quelle1,Quelle2,...,QuelleN}|
explizit referenziert werden.

Sofern man eine einzige Datenbank für sämtliche Literaturverzeichnisse verwendet
und die unsichtbare Referenzierung jeweils mit dem \lstinline|\nocite{*}|"=Befehl erfolgt,
müssen die relevanten Referenzen für das jeweilige Zusatz"=Literaturverzeichnis
gefiltert werden.
Die Filterung der eigenen Publikationen, Patente und betreuten Arbeiten
für die einzelnen Verzeichnisse geschieht jeweils mit der Angabe
\lstinline|keyword=ownpubl|,
\lstinline|keyword=ownpatent|
oder
\lstinline|keyword=supervised|
in der Optionenliste des jeweiligen
\lstinline|printbibliography|"=Befehls.
Dazu müssen die Literatureinträge in der Literaturdatenbank freilich
entsprechend jeweils mit der Angabe
\lstinline|keywords = {ownpubl}|,
\lstinline|keywords = {ownpatent}|
oder
\lstinline|keywords = {supervised}|
versehen werden.

Bei Verwendung separater Literaturdatenbanken für die einzelnen Zusatzverzeichnisse
oder bei einer expliziten Aufzählung der einzelnen Veröffentlichungen mit dem
\lstinline|\nocite{Quelle1,Quelle2,...,QuelleN}|"=Befehl,
ist eine Filterung der Literaturdatenbank nach einem Keyword nicht notwendig.
Damit könnte die Notwendigkeit von Zusatzangaben in der Literaturdatenbank entfallen.
Allerdings muss dann auch die standardmäßig gesetzte Filter-Option in den jeweiligen
\lstinline|printbibliography|"=Befehlen
in der Datei  \printfilepathquoted{content/Inhalt-BackMatter.tex}
deaktiviert (sprich auskommentiert) werden.

Die Zusatzauflistungen sind von den Zitierungen im Text unabhängig.
Damit ergibt sich für eigene Veröffentlichungen,
die im Text explizit zitiert werden, eine doppelte Listung --
einmal in der Hauptliteraturliste (mit einer alphanumerischen Referenz)
und einmal in der durchnummerierten Liste eigener Publikationen.
Dies ist kein Bug, sondern ein Feature der Vorlage.

Bei der Auflistung eigener Veröffentlichungen und Patente ist 
in der Vorlage jeweils ein kleines Vorwort vorgesehen.
Darin können beispielsweise Zusatzhinweise zur Abgrenzung der Publikationen \bzw
der Patente, die während der Promotionszeit entstanden sind,
oder Hinweise zur Eingruppierung der Veröffentlichungen platziert werden.
Diese einleitenden Texte sind ebenfalls in der Datei
\printfilepathquoted{content/Inhalt-BackMatter.tex}
definiert und können dort editiert werden.

Bei der Erwähnung der einzelnen eigenen Veröffentlichungen und Patente in den Vorworten 
zu den jeweiligen Listen
muss die Einbindung mit Hilfe des Makros
\lstinline|\citefield{...}[labelnumberwidth]{labelnumber}|
statt der üblichen Einbindung mit dem Makro
\lstinline|\cite{...}|
erfolgen, damit die zu den Listen passende (nummerische) Formatierung verwendet wird.

Sind keine Vorworte vor den Listen eigener Veröffentlichungen und Patente erwünscht,
lassen sich diese durch die Auskommentierung der Optionen
\lstinline|prenote=OwnPublPrenote|
\bzw
\lstinline|prenote=OwnPatentPrenote|
in den jeweiligen 
\lstinline|printbibliography|"=Befehlen komplett ausblenden.
%%%%%%%%%%%%%%%%%%%%%%%%%%%%%%%%%%%%%%%%%%%%%%%%%%%%%%%%%%%%
\section{Abkürzungsverzeichnis, Stichwortverzeichnis (Index), Symbolverzeichnis und Glossar}%
\index{Abkürzungsverzeichnis}%
\index{Glossar}%
\index{Symbolverzeichnis}%
\index{Stichwortverzeichnis}%
%\index{Index|see{Stichwortverzeichnis}}
\label{sec:Glossare}
%%%%%%%%%%%%%%%%%%%%%%%%%%%%%%%%%%%%%%%%%%%%%%%%%%%%%%%%%%%%
%
Die Vorlage unterstützt auch ein Abkürzungsverzeichnis, ein Stichwortverzeichnis,
ein Symbolverzeichnis sowie ein allgemeines \index{Glossar}\gls{gls:Glossar},
das Definitionen von \index{Fachbegriff!Definition}Fachtermini oder
\index{Uebersetzung@Übersetzung}Übersetzungen von fremdsprachlichen Begriffen
(d.h. eine Art Lexikons) enthalten kann.


%%%%%%%%%%%%%%%%%%%%%%%%%%%%%%%%%%%%%%%%%%%%%%%%%%%%%%%%%%%%
\subsection*{Abkürzungen und Abkürzungsverzeichnis}%
%% Da kein Eintrag für dieses Kapitel im Imhaltsverzeichnis:
%% PDF-Lesezeichen setzen:
\pdfbookmark[2]{Abkürzungen und Abkürzungsverzeichnis}{abk}
\index{Abkürzungen}%
\index{Akronyme}%
\label{sec:Akronyme}
%%%%%%%%%%%%%%%%%%%%%%%%%%%%%%%%%%%%%%%%%%%%%%%%%%%%%%%%%%%%
Zur Erzeugung des Abkürzungsverzeichnises und des Glossars wird intern das
\pkg{glossaries}-\gls{gls:pkg} verwendet \cite{talbot2014}.
Zudem wurden einige Makros definiert, welche die Erfassung der Begriffe erleichtern sollen.
%
Um ein Abkürzungsverzeichnis zu erzeugen, muss zuerst eine Liste der Abkürzungen angelegt werden.
Dies geschieht in der Datei \printfilepathquoted{./preambel/Acronyms.tex}.
Hierfür wird das Makro \lstinline|\newacronym| verwendet.
Dieses Makro hat drei obligatorische Argumente, nämlich
\begin{itemize*}
\item die Marke,
\item das Akronym und
\item die Langform.
\end{itemize*}

Als Konvention wird der Marke eines Akronyms \printkeyword{ac:} als Präfix vorangestellt.

\begin{latex}[caption={Definition einer Abkürzung},label={lst:SimpleAcronymEntry}]
\newacronym{ac:MSA}{MSA}{mein schönes Akronym}
\end{latex}

Optional können Pluralformen, sowie Genitiv-, Dativ- und Akkusativ-Formen 
der Abkürzung und des eigentlichen Begriffs angegeben werden,
sofern sie im Quelltext verwendet werden und sich von der Grundform unterscheiden.
Außerdem kann mit dem Schlüsselwort \printkeyword{description}
eine abweichende Version der Langform für die Verwendung im Abkürzungsverzeichnis definiert werden:

\begin{latex}[caption={Definition einer Abkürzung mit Zusatzangaben},label={lst:AdvancedAcronymEntry}]
\newacronym[shortgenitive={MSAs},%
            genitive={meines schönen Akronyms},%
            %shortdative={MSA},% <- entspriht der Hauptform 
            %            % ==> keine Zusatzdefinition nötig!
            dative={meinem schönen Akronym},%
            %shortaccusative={MSA},%
            %accusative={mein schönes Akronym},%
            shortplural={MSAs},
            longplural={meine schönen Akronyme},%
            %shortpluralgenitive={MSAs},%
            pluralgenitive={meiner schönen Akronyme},%
            %shortpluraldative={MSAs},%
            pluraldative={meinen schönen Akronymen},%
            %shortpluralaccusative={MSAs},%
            pluralaccusative={meine schönen Akronyme},%
            description={mein schönes Akronym, % <- optional!
                         ein Beispiel für eine Abkürzung}%
           ]{ac:MSA}{MSA}{mein schönes Akronym}
\end{latex}

Im Text des Dokumentes werden die Einträge durch den Befehl \lstinline|\ac{<Marke>}| verwendet.
Bei der erstmaliger Verwendung wird die Langform gedruckt, gefolgt von der Abkürzung, welche in Klammern gesetzt wird.
Beim zweiten Vorkommen wird nur noch die Abkürzung gedruckt.
Zusätzlich definiert die Vorlage die Befehle
\lstinline|\acgen{...}|, \lstinline|\acdat{...}| und \lstinline|\acacc{...}|, sowie
\lstinline|\acplgen{...}|, \lstinline|\acpldat{...}| und \lstinline|\acplacc{...}|,
die jeweils die Genitiv-, Dativ- und Akkusativ-Form (Singular und Plural) drucken.

Im \gls{gls:pkg} \pkg{glossaries} stellt der Befehl \lstinline|\ac{...}| ein Shortcut für den Befehl \lstinline|\gls{...}| dar.
Mit diesem kann ein allgemeiner Glossar-Eintrag im Text referenziert werden.

Mit den Befehlen
\lstinline|\acrshort{...}|, \lstinline|\acrlong{...}|, \lstinline|\acrfull{...}|
und ihren Abwandlungen sowie den Shortcuts
\lstinline|\acs{...}|, \lstinline|\acl{...}|, \lstinline|\acf{...}|,
kann jeweils nur die Abkürzung, nur die Langform oder beides explizit angefordert werden.
Allerdings wird eine solche Verwendung ggf. nicht als \enquote{erstmalige Verwendung} zählen.

Damit resultiert der folgende Quellcode
\begin{latex}[caption={Verwendung von Abkürzungen},label={lst:AcronymUsage}]
    \Acf{ac:MSA} ist ein Beispiel für die Verwendung
    einer Abkürzung am Anfang des Satzes. Man beachte,
    dass der Aufruf der Marke mit dem Makro
    \lstinline|\acf{...}| bzw. \lstinline|\Acf{...}|
    nicht als die erste Erwähnung im Text zählt.
    Bei der ersten Nennung \acgen{ac:MSA} unter
    Verwendung der Makros \lc{ac}, \lc{acgen} \oae
    erscheint die Langform, gefolgt von der Kurzform.
    Bei der zweiten Nennung \acgen{ac:MSA}
    erscheint nur noch die Kurzform \ac{ac:MSA}.
\end{latex}
%
in der folgenden Ausgabe:
%
\begin{quote}
    \Acf{ac:MSA} ist ein Beispiel für die Verwendung
    einer Abkürzung am Anfang des Satzes. Man beachte,
    dass der Aufruf der Marke mit dem Makro
    \lstinline|\acf{...}| bzw. \lstinline|\Acf{...}|
    nicht als die erste Erwähnung im Text zählt.
    Bei der ersten Nennung \acgen{ac:MSA} unter
    Verwendung der Makros \lc{ac}, \lc{acgen} \oae
    erscheint die Langform, gefolgt von der Kurzform.
    Bei der zweiten Nennung \acgen{ac:MSA}
    erscheint nur noch die Kurzform \ac{ac:MSA}.
\end{quote}


%%%%%%%%%%%%%%%%%%%%%%%%%%%%%%%%%%%%%%%%%%%%%%%%%%%%%%%%%%%%
\subsection*{Glossar}%
%% Da kein Eintrag für dieses Kapitel im Imhaltsverzeichnis:
%% PDF-Lesezeichen setzen:
\pdfbookmark[2]{Glossar}{gloss}
\index{Glossar}%
\label{sec:Glossar}
%%%%%%%%%%%%%%%%%%%%%%%%%%%%%%%%%%%%%%%%%%%%%%%%%%%%%%%%%%%%

Neben einem Abkürzungsverzeichnis kann man auch ein
\gls{gls:Glossar} erstellen lassen.
In diesem können Definitionen von \index{Fachbegriff!Definition}Fachtermini oder
\index{Uebersetzung@Übersetzung}Übersetzungen von fremdsprachlichen Begriffen stehen.

Der wesentliche Unterschied zwischen einer Abkürzung und einem allgemeinen Glossar-Eintrag ist,
dass bei Abkürzungen bei erstmaliger Verwendung die Abkürzung gedruckt und die Langform in Klammer dahinter gesetzt wird.
Bei allgemeinen Glossar-Einträgen wird normalerweise nur der Name gesetzt.
Durch eine Option des \pkg{glossaries}-\glsgen{gls:pkg} kann man sicher stellen,
dass alle Glossar-Einträge auf das Glossar am Ende des Manuskripts verlinkt werden.
Standardmäßig ist diese Option jedoch deaktiviert.

Die Glossar-Einträge werden in der Datei  \printfilepathquoted{./preambel/Glossary.tex} definiert.

Die Definition der Glossar-Einträge geschieht mit Hilfe des eigens für diese Vorlage definierten Makros \lc{myglossaryentry},
welches drei obligatorische Argumente hat, nämlich
\begin{itemize*}
\item die Marke,
\item den Begriff und
\item die Erklärung / Definition / Übersetzung.
\end{itemize*}
Als Konvention wird der Marke eines Glossar-Eintrages ein \enquote{gls:} als Präfix vorangestellt.
Optional können die Genitiv-, Dativ-, Akkusativ-" sowie die Plural=Formen des Begriffs angegeben werden,
sofern sie im Quelltext verwendet werden und sich von der Grundform unterscheiden.

Ein Glossar-Eintrag kann beispielsweise folgendermaßen definiert werden:
\begin{latex}[caption={Definition eines Glossar-Eintrages},label={lst:GlossEntry}]
\myglossaryentry[genitive={Glossars}, plural={Glossare},
                 pluraldative={Glossaren}]%
                {gls:Glossar}{Glossar}{alphabetisch sortierte Liste von Begriffen mit Erklärung}
\end{latex}

Normalerweise werden im Glossar nur diejenigen Begriffe angezeigt,
die im Text des Dokumentes erwähnt und entsprechend referenziert worden sind.
Eine Referenzierung der Glossar-Einträge im Text geschieht normalerweise mit dem
\lstinline|\gls{<Marke>}|-Befehl,
welcher den Begriff im Text druckt und für seine Aufnahme ins Glossar sorgt.
Weitere mögliche Befehle sind \lstinline|\Gls{...}| und \lstinline|\GLS{...}|,
die den ersten \bzw alle Buchstaben in Großbuchstaben umwandeln,
\lstinline|\glspl{...}|, \lstinline|\Glspl{...}|, \lstinline|\GLSpl{...}|
für die Pluralform usw.
Zusätzlich definiert die Vorlage die Befehle
\lstinline|\glsgen{...}|, \lstinline|\glsdat{...}| und \lstinline|\glsacc{...}|,
sowie
\lstinline|\glsplgen{...}|, \lstinline|\glspldat{...}| und \lstinline|\glsplacc{...}|,
die jeweils die Genitiv-, Dativ- und Akkusativ-Form drucken.
So kann beispielsweise durch die Verwendung des Befehls 
\lstinline|\glsgen{gls:Glossar}| eine Erwähnung des oben definierten \enquote{Glossar}-Eintages
in der Genitiv-Form bewirkt werden, also eine Erwähnung des \glsgen{gls:Glossar}.

Mit dem Befehl \lstinline|\glsadd{...}| gibt es zusätzlich die Möglichkeit,
eine Stelle im Text mit einem Glossar-Begriff zu verlinken, ohne diesen explizit zu drucken.
Mit \lstinline|\glsaddall| kann man alle definierte Glossar-Einträge ins Glossar aufnehmen,
ohne sie im Text des Dokumentes referenziert zu haben.

Ein Glossar-Eintrag kann bei Bedarf mit \glspldat{gls:Untereintrag} versehen werden.
Diese können wie folgt definiert werden:
\begin{latex}[caption={Definition eines Glossar-Untereintrages},label={lst:GlossSubentry}]
\myglossarysubentry[genitive={Untereintrags},
      plural={Untereinträge},pluraldative={Untereinträgen}]%
     {gls:Untereintrag}{gls:Glossar}{Untereintrag}{ein untergeordneter Eintrag im Glossar}
\end{latex}
Eine Untereintrag-Definition hat vier obligatorische Argumente, nämlich
\begin{itemize*}
\item die Marke,
\item die Marke des übergeordneten Eintrags,
\item den Begriff und
\item die Erklärung / Definition / Übersetzung.
\end{itemize*}


%%%%%%%%%%%%%%%%%%%%%%%%%%%%%%%%%%%%%%%%%%%%%%%%%%%%%%%%%%%%
\subsection*{Symbolverzeichnis}%
%% Da kein Eintrag für dieses Kapitel im Imhaltsverzeichnis:
%% PDF-Lesezeichen setzen:
\pdfbookmark[2]{Symbolverzeichnis}{symbolverz}
\index{Glossar}%
\index{Symbolverzeichnis}%
\index{Symbole}%
\index{Notation}%
\label{sec:Symbolverz}
%%%%%%%%%%%%%%%%%%%%%%%%%%%%%%%%%%%%%%%%%%%%%%%%%%%%%%%%%%%%
%
Ein Symbolverzeichnis kann auf zwei Arten angelegt werden.
Normalerweise reicht eine manuell erstellte Übersicht über die Notation,
so wie sie in der Datei \printfilepathquoted{./00-Front-Matter/Notation.tex}
definiert wurde. Hierzu wurden in der Vorlage zwei Arten von Tabellenumgebungen
\printkeyword{myNotationTable} und \printkeyword{myNotationDescTable}
(zwei- \bzw dreispaltig) mit unsichtbaren Rändern sowie folgende Makros
für die entsprechenden Zeileneinträge definiert.

Für die zweispaltige Tabelle:
\begin{itemize*}
\item \lstinline|\myNotationTableEntryMath{<Mathe-Ausdruck>}{<Text-Beschreibung>}| und
\item \lstinline|\myNotationTableEntryText{<Text>}{<Text>}|
\end{itemize*}
und für die dreispaltige Tabelle:
\begin{itemize*}
\item \lstinline|\myNotationDescTableEntry{<Text>}{<weiterer Text>}{<Mathe-Ausdruck>}|,
\item \lstinline|\myNotationDescTableEntryText}{<Text>}{<Text>}{<Text>}|.
\end{itemize*}

Allerdings gibt es auch die Möglichkeit zur automatischen Erzeugung eines
Symbolverzeichnisses mit Hilfe des \pkg{glossaries}-\glsgen{gls:pkg}.
Um dieses automatisch generierte Symbolverzeichnis zu erzeugen,
müssen Symbole in Form von Glossar-Einträgen angelegt und im Text des Dokumentes
zumindest einmal entsprechend mit dem Befehl \lstinline|\gls{<Marke>}| referenziert werden.
Der Eintrag für das Symbol \gls{symb:pi} müsste beispielsweise folgendermaßen angelegt werden:
%\newpage
\begingroup%
\setlength{\topskip}{0pt}%
\begin{latex}[caption={Definition eines Symboleintrages},label={lst:SymbEntry}]
\newglossaryentry{symb:pi}{%
                  name={\ensuremath{\pi}},%
                  sort={pi},%
                  type=symbols,%
                  description={Kreiszahl, Verhältnis des Umfangs eines Kreises zu seinem Durchmesser}%
                 }
\end{latex}
\endgroup%

In der Vorlage sind beide Arten von Symbolverzeichnissen vorgesehen.
Sie werden am Anfang des Manuskriptes vor dem Inhaltsverzeichnis eingebunden.
Es passiert in der Datei \printfilepathquoted{./content/Inhalt-FrontMatter.tex}.
Bei Bedarf lässt sich diese Einbindung auch an eine geeignete Stelle in den Backmatter verschieben.

Zum einfachen Ein- \bzw Ausblenden der beiden Symbolverzeichnisse
gibt es in der Datei \printfilepathquoted{preambel/AlleSchalter.tex}
die beiden Schalter
\begin{itemize*}
\item \printkeyword{showNotation} für den manuell erstelle Notationübersicht
\item \printkeyword{showListOfSymbols} für das automatisch erstelltes Symbolverzeichnis.
\end{itemize*}



%%%%%%%%%%%%%%%%%%%%%%%%%%%%%%%%%%%%%%%%%%%%%%%%%%%%%%%%%%%%
\subsection*{Stichwortverzeichnis (Index)}%
%% Da kein Eintrag für dieses Kapitel im Imhaltsverzeichnis:
%% PDF-Lesezeichen setzen:
\pdfbookmark[2]{Stichwortverzeichnis}{stichwortverz}
\index{Stichwortverzeichnis}%
\index{Index|see{Stichwortverzeichnis}}%
\label{sec:Index}
%%%%%%%%%%%%%%%%%%%%%%%%%%%%%%%%%%%%%%%%%%%%%%%%%%%%%%%%%%%%
%
Ein \index{Stichwortverzeichnis}Stichwortverzeichnis (oder Index)
ist eine alphabetisch sortierte Liste von Begriffen mit einer Auflistung der Fundstellen im Dokument.
Diese ist nützlich, wenn sich der Leser zu einem Begriff alle Vorkommnisse sehen möchte.
Der Index wird erzeugt, indem im Quellcode der Befehl \lstinline+\index{Begriff}+
eingefügt wird. Der Begriff selbst wird dadurch nicht gedruckt und muss daher
noch einmal wiederholt werden, um auch im Text gedruckt zu werden. Dieses Verhalten
ist beabsichtigt, sodass im Index immer nur die Grundform des Wortes verwendet
wird, aber im Text natürlich die richtige Deklination.

%%%%%%%%%%%%%%%%%%%%%%%%%%%%%%%%%%%%%%%%%%%%%%%%%%%%%%%%%%%%
\section{Randnotizen und TODO-Notizen}%
\index{Randnotizen}%
\index{TODO-Notizen}%
\label{sec:Randnotizen}
%%%%%%%%%%%%%%%%%%%%%%%%%%%%%%%%%%%%%%%%%%%%%%%%%%%%%%%%%%%%
%
Randnotizen 
\floatmarginnote{Ich bin eine überflüssige Randnotiz}%
werden mit dem Kommando \lc{floatmarginnote} gesetzt.
Sie eignen sich zum Beispiel, um wichtige Begriffe oder Aussagen zu verdeutlichen
oder durch eine Kurzzusammenfassung jedes einzelnen Textabschnitts
den roten Faden zu verdeutlichen.
Aufgrund des sehr kleinen Randbereichs im Seitenlayout
können in der endgültigen Fassung des Manuskriptes keine Randnotizen platziert werden.
Um \ggf vorhandene Randnotizen auszublenden,
muss man in der Datei \printfilepathquoted{preambel/AlleSchalter.tex} den Wert des Schalters
\printkeyword{showMarginNotes} auf \printkeyword{false} setzen.

TODOs im Text lassen sich mit Hilfe des Befehls
\todo{Hier weitere Beispiele hinzufügen!}
\lc{todo\{<Hinweistext>\}} aus dem Paket \pkg{todonotes} setzen.
Dies geschieht wahlweise als Randnotiz
oder als eine Box mitten im Text mit der Option \printkeyword{inline}:
\lc{todo[inline]\{<Hinweistext>\}}.
\todo[inline]{Inline-TODO: Hier noch mehr Beispiele hinzufügen!}

Ähnlich erzeugt der Befehl \lc{missingfigure\{<Hinweistext>\}}
einen Platzhalter für eine noch zu erstellende Abbildung.
\missingfigure{Hier fehlt eine Grafik!}

Die TODOs lassen sich bei Bedarf ausblenden, indem man in der Datei
\printfilepathquoted{preambel/AlleSchalter.tex} den Schalter
\printkeyword{showTODOs} auf \printkeyword{false} setzt.

Sofern TODOs vorhanden und eingeblendet sind, wird am Ende des Manuskriptes automatisch eine \hyperlink{TODO-List}{Übersicht der TODOs} engefügt.
%%%%%%%%%%%%%%%%%%%%%%%%%%%%%%%%%%%%%%%%%%%%%%%%%%%%%%%%%%%%
\section[Anpassungen für den Druck]{Anpassungen für den Druck}%
\index{Druck}%
\label{sec:Druckanpassungen}
%%%%%%%%%%%%%%%%%%%%%%%%%%%%%%%%%%%%%%%%%%%%%%%%%%%%%%%%%%%%
%
Vor der Abgabe einer Druckversion an den KSP Verlag sollte die korrekte Schaltersetzung 
in der Datei \printfilepathquoted{./preambel/AlleSchalter.tex} nochmals geprüft werden.
Insbesondere sollten mit der temporären Umstellung des Schalters
\printkeyword{showFrame} auf \printkeyword{true}
die Satzspiegel"=Ränder testweise angezeigt werden und so kontrolliert werden,
ob nichts hinausragt und ob Tabellen und Bilder die komplett verfügbare Breite ausfüllen.
Besonders sollte man dabei auch auf die Overfullbox- und Underfullbox"=Warnungen achten und diese beheben.
Eine Overfullbox"=Warnung tritt dann auf, wenn ein Objekt (Text, Bild, Tabelle)
über den dafür vorgesehenen Bereich (normalerweise über den Satzspiegelrand) hinausragt.
Eine Underfullbox"=Warnung tritt auf, wenn durch fehlende Möglichkeit der Zeilenumbrüche der Text zu stark gedehnt werden müsste und die Wortzwischenräume zu groß sind.
Hierzu sollte man insbesondere auch Hinweise zur Silbentrennung aus
\cref{sec:Hyphenation,sec:Titles} beachten.
Danach sollte der Schalter \printkeyword{showFrame} wieder auf \printkeyword{false} umgestellt werden.

Bei Abbildungen, die am Anfang einer Seite stehen, ist darauf zu achten,
dass es keinen leeren Zwischenraum zwischen dem oberen Rand des Satzspiegels und der Abbildung gibt.
Normalerweise passiert dies durch die Einstellungen der Vorlage automatisch.
\index{Bild!zuschneiden}%
\index{Bild!weißer Rand}%
\index{Bild!clip}%
\index{Bild!trim}%
Falls doch noch ein leerer Raum entsteht, liegt dies zumeist daran, dass die Abbildung selbst weiße Ränder hat.
Bei Binärbildern können diese durch Setzen der Optionen
\printkeyword{clip=true} und \printkeyword{trim=<left> <lower> <right> <upper>}
im \printkeyword{includegraphics}-Befehl zugeschnitten werden, \zB so:
\lstinline|\includegraphics[clip=true, trim={0cm 1cm 0cm 1cm}]{Bildpfad\Dateiname}|.

Mit den beiden Einstellungen 
\lstinline|\setUserDefinedBoolean{coloredlistings}{false}| und 
\lstinline|\setUserDefinedBoolean{coloredlinks}{false}|
sollte dafür gesorgt werden, dass alle Quellcode-Listings sowie Querverweise und URL-Adressen,
die normalerweise farbig sind, für die Druckversion nicht farbig gesetzt werden
und so die Anzahl der farbig zu druckenden Seiten reduziert wird.

Mit der Einstellung
\lstinline|\setUserDefinedBoolean{useCMYKcolors}{true}|
eine Farbkonvertierung aller Farben in den CMYK-Farbraum für den Offset-Druck vorgenommen werden.
So kann man die Farbgebung der Druckversion begutachten und sicherstellen, dass alle Farben beim Druck gut aussehen.%
\footnote{Da manche RGB-Farben bei der Konvertierung in den CMYK-Farbraum blass aussehen,
ist es sinnvoll, eine Alternativversion der Farben im CYMK-Farbraum zu definieren,
die gut aussieht.
Beispiele dafür gibt es in der Datei \printfilepathquoted{./preambel/ColorSettings.tex}.
Vorzugsweise sollten aber die KIT-Corporate-Identity-Farben verwendet werden,
welche in der Datei \printfilepathquoted{KAcolors.sty} definiert sind.
Für diese Farben wurde sowohl eine RGB- als auch eine CMYK-Definition erstellt.}


Auch gibt es zwei Sachen, die in der Datei 
\printfilepathquoted{./preambel/AlleAngaben.tex}
geändert werden müssen.
Zum einem ist der Status der Dissertation von \enquote{vorgelegt} auf \enquote{genehmigt} zu ändern
(dies geschieht durch die Umstellung des Schalters
\printkeyword{statusApproved}
von \printkeyword{false} auf \printkeyword{true}).
Sofern erst nach der Prüfung entschieden wird, ob man beispielsweise Dr.-Ing.\ oder Dr.\ rer.\ nat.\ wird, 
ist die \printkeyword{DocDegree}-Angabe entsprechend anzupassen.

\todo[inline,caption={Vor der Abgabe alle TODOs und die Übersichtsliste ausblenden!}]{%
Vor der Abgabe sollte man sich vergewissern, dass alle TODOs erledigt sind
und die Liste der noch zu erledigenden Punkte am Ende des Manuskriptes nicht mehr erscheint.
Die ggf. noch verbleibenden TODOs und die Auflistung dieser am Ende des Manuskriptes lassen sich ausblenden,
indem man in der Datei \printfilepathquoted{./preambel/AlleSchalter.tex}
den Schalter \printkeyword{showTODOs} auf \printkeyword{false} setzt.
}
%
\end{showExamples}%